\documentclass[12pt, a4paper]{academics}

\academicsCoverLabel{解题报告}
\academicsCourse{算法设计与分析}{Design and Analysis of Algorithms}
\academicsReport{蓝桥杯 2026 国赛 A 组题解}{2026 年 6 月 11 日}
\academicsArthur{华博文}{19240212}

\begin{document}

\makeacademicscover

\section{引言}

“蓝桥杯”全国软件和信息技术专业人才大赛是由工业和信息化部人才交流中心主办的全国性 IT 学科赛事，至今已连续举办十余届，累计参赛人数超过百万。2026 年蓝桥杯 C/C++ 程序设计竞赛国赛 A 组共包含 10 道题目，涵盖数论推导、组合计数、动态规划、图论模型、博弈论、贪心算法与数据结构等多个算法领域，对选手的数学建模能力与代码实现效率提出了较高要求。

本报告针对该 10 道题目逐一给出完整的题解。从问题重述出发，阐明数学模型与核心推导，分析时间复杂度与空间复杂度，最后附以可运行的 C++ 代码实现。题目按编号依次为：密钥提取（01）、圆形灯环（02）、多项式积分（03）、安全路径（04）、生态廊道（05）、独立三角形（06）、核心任务调度（07）、魔法前缀（08）、神秘排列（09）和不互质游戏（10）。其中题 02 涉及 Burnside 引理与矩阵快速幂，题 08 与题 10 涉及 Sprague--Grundy 定理与公平组合游戏分析，题 06 与题 07 涉及贪心策略与拟阵结构，题 09 涉及错位排列的容斥计数与大规模模运算。

报告中的所有代码均使用 C++23 标准编写，在 Apple clang 21.0.0 编译器下通过测试。

\section{题01：密钥提取（P16783）}

\subsection{从重复数字到代数恒等式}

密钥$A$由$n = 20260606$个数字$9$拼接而成，密钥$B$由$m = 2026$个数字$9$拼接而成。在十进制中，$k$个连续的$9$具有简洁的代数形式：$\underbrace{99\ldots9}_{k\text{个}} = 10^{k} - 1$。取$k=3$可得$999 = 10^3-1$，取$k=5$可得$99999 = 10^5-1$，这一规律对任意$k$均成立。据此，$A = 10^n - 1$，$B = 10^m - 1$，两者的乘积展开为
\[
S = (10^n - 1)(10^m - 1) = 10^{n+m} - 10^n - 10^m + 1.
\]
题目要求提取$S$的十进制表示中（不含前导零，从左往右以$1$为起始编号）第$2026$位、第$20260523$位和第$20262632$位的三个数字$d_1, d_2, d_3$，并以它们组成的三位数$\overline{d_1 d_2 d_3}$作为答案。至此，问题从“两个巨大全$9$数的乘积”转化为了“四个$10$的幂的加减结果中特定位置的数字”，后者可以通过分析代数结构直接求解。

\subsection{暴力乘法的不可行性}

如果试图直接计算$A \times B$，$A$本身约$2 \times 10^7$位十进制数字，存储这个数就需要约$20$ MB内存，乘积$S$更是拥有$n+m = 20262632$位。即便采用基于快速傅里叶变换的大整数乘法，时间复杂度也在$\mathrm{O}(N \log N)$级别，其中$N \approx 2 \times 10^7$为乘积位数，运算量接近$5 \times 10^8$次浮点操作，远超竞赛时限和内存约束。况且这是一道结果填空题，题目设计者显然期望选手绕过数值计算、从数的结构入手。因此放弃暴力乘法、转向代数模式分析是唯一可行的路径。

\subsection{交换减法顺序：一个简洁的推导}

表达式$S = 10^{n+m} - 10^n - 10^m + 1$包含四个项，可以通过不同的分组顺序来化简。一种直观的思路是先合并$10^{n+m} - 10^n$，但该减法涉及跨越多位的借位链，手工追踪较为繁琐。更巧妙的方式是先合并后两项的相反顺序：将$10^{n+m}$与$-10^m$组合。注意到
\[
10^{n+m} - 10^m = 10^m(10^n - 1),
\]
而$10^n - 1$恰好是$n$个$9$，乘以$10^m$相当于在末尾追加$m$个$0$。因此$10^{n+m} - 10^m$的十进制表示极为规则：最高$n$位全部是$9$，紧随其后的$m$位全部是$0$。以小规模参数为例：取$n = 5$，$m = 3$，则$10^3 \times (10^5 - 1) = 10^3 \times 99999 = 99999000$，恰好是$5$个$9$后接$3$个$0$，共$8$位，最高位是$9$，最低位是$0$。

接下来从这个“$n$个$9$、$m$个$0$”的数中减去$10^n$。$10^n$在一个$n+m$位的大数中对齐到$10^n$的幂次位置（从右数第$n+1$位，即从左数第$m$位），该位在上述数中的原值是$9$（因为它落在前$n$位的$9$区域之内——这里依赖$m < n$，本题中$2026 \ll 20260606$显然满足）。由于被减位本就是$9$而非$0$，减法极其简单：$9 - 1 = 8$，且不产生任何向前借位。减完后，第$m$位从$9$变为$8$，其余位不变。最后加上末尾的$+1$：最低位（第$n+m$位）从$0$变为$1$，且因为$n+m-1$位也是$0$，$1$加到$0$上不产生进位。

仍以$n=5, m=3$为例走完整条推导链。第一步，$10^3(10^5-1) = 99999000$（位置$1\sim 5$为$9$，$6\sim 8$为$0$）。第二步，减去$10^5 = 100000$：$10^5$的$1$落在第$3$位（因为总位数$8$，$10^5$对应幂次$5$，从左数位置$8-5=3$），该位置原为$9$，减$1$得$8$；其余位不受影响。结果变为$99899000$。第三步，加$1$：末位$0$变为$1$，得$S = 99899001$。逐位罗列这些数字：$9, 9, 8, 9, 9, 0, 0, 1$。由此可以归纳出$S$在任何一般参数下的完整数字布局。

\subsection{六大区间：积的完整数字模式}

将上述推导推广到一般的$n$和$m$（设$m < n$，本题天然满足），$S$的$n+m$位十进制数字从高位到低位可划分为六个连续区间。第一区间，第$1$位至第$m-1$位，全部是$9$（减法中唯一改变的是第$m$位，这些更高位的$9$完全未受影响）。第二区间，仅含第$m$位，其值为$8$（被$10^n$减走了$1$）。第三区间，第$m+1$位至第$n$位，全部是$9$（同属前$n$个$9$的块，但位于第$m$位右侧，同样未受减法影响）。第四区间，仅含第$n+1$位，其值为$0$（这是后$m$个$0$的开头，$+1$的进位无法到达这里）。第五区间，第$n+2$位至第$n+m-1$位，全部是$0$（同为后$m$个$0$的延续）。第六区间，仅含最末位第$n+m$位，其值为$1$（来自末尾的$+1$）。

为增强信心，用另一组参数交叉验证：取$n = 7$，$m = 4$（总位数$11$）。按相同步骤，$10^4(10^7-1) = 99999990000$，减$10^7$（落在第$4$位）得$99989990000$，加$1$得$S = 99989990001$。观察其$11$位：位置$1\sim 3$为$9$、$4$为$8$、$5\sim 7$为$9$、$8$为$0$、$9\sim 10$为$0$、$11$为$1$，与上述六个区间的划分完全吻合。

\subsection{定位三位目标数字}

将题目参数$n = 20260606$，$m = 2026$代入，总位数$n+m = 20262632$。首先确定第一位$d_1$，位于第$2026$位。该位置恰好等于$m$，落入第二区间（唯一的$8$所在位），故$d_1 = 8$。其次确定第二位$d_2$，位于第$20260523$位。该位置满足$m < 20260523 \le n$（具体而言$20260523 = n - 83$），落入第三区间（$9$的区域），故$d_2 = 9$。最后确定第三位$d_3$，位于第$20262632$位。该位置恰好等于总位数$n+m$，落入第六区间（末尾的$1$），故$d_3 = 1$。因此所求三位整数为$\overline{d_1 d_2 d_3} = \overline{891} = 891$。

\subsection{常数时间的优雅解法}

整个求解过程不涉及任何大整数运算，仅使用整数$n$和$m$的加减与大小比较——将两个全$9$数识别为$10^k - 1$、展开乘积后交换减法顺序以避免借位、归纳六个数字区间、判定三个给定位置各自落入哪个区间。每一步均为$\mathrm{O}(1)$的整数操作，空间开销也仅需存储几个整型变量，同为$\mathrm{O}(1)$。将一个表面上看需要处理约$2 \times 10^7$位数字的“天文”问题化为常数时间，核心在于两条洞察：第一，$k$个$9$恒等于$10^k - 1$，将大数乘法转化为$10$的幂的加减；第二，交换减法顺序（先做$10^{n+m}-10^m$再做$-10^n$而非先做$-10^n$），恰好使被减位落在$9$上从而避免复杂的借位链。这种“用代数结构代替数值计算”的策略，正是竞赛中面对超大数字问题的标准思路。

\subsection{代码实现}
\inputminted{c++}{../src/p01_key_extraction.cpp}

\section{题02：圆形灯环（P16784）}

\subsection{问题建模：灯环约束转化为有向图上的闭游走}

一个由$N=2026$个LED灯珠首尾相连组成的圆形灯环，每个灯珠的亮度可在$0\sim 25$共$K=26$个等级中独立选择。一个配置“稳定”需满足两个条件：顺时针方向上任意连续$3$个灯珠的亮度之和$\neq 26$；相邻对$(2,6)$（前一个为$2$级、后一个为$6$级）在整个环中出现总次数为偶数。能通过顺时针旋转相互重合的配置视为本质相同。要求计算本质不同的稳定配置数模$998244353$的结果。

核心思路是将环形序列建模为有向图上的闭游走。定义状态为相邻两灯珠的亮度对$(x,y)$，共有$K^2=26^2=676$个状态，每个状态用编号$x\cdot K+y$唯一标识。状态$(x,y)$可一步转移到$(y,z)$当且仅当三元组约束不触犯，即$x+y+z\neq 26$。例如状态$(5,10)$：枚举$z=0\sim 25$，$5+10+z=26$解得$z=11$，故$z=11$被禁止，其余$25$个$z$值均可转移；状态$(0,0)$：$0+0+z=26$要求$z=26$超出范围，因此全部$26$个$z$均可转移；状态$(12,14)$：$12+14+z=26$给出$z=0$，仅$z=0$被禁止。每个状态的出度在$25$或$26$之间。

考虑一个长度为$d$的环形配置$a_0,a_1,\dots,a_{d-1}$（下标模$d$）。以$(a_0,a_1)$为起点，依次转移至$(a_1,a_2),\dots,(a_{d-1},a_0)$，最终回到$(a_0,a_1)$，这恰好是一条长度为$d$的闭游走。游走的$d$次转移依次检验了环上全部$d$个三元组约束，且访问的$d$个状态正是环上$d$个相邻对。关键观察：从形序列到闭游走的映射是一一且映上的——给定闭游走可唯一还原环形序列$a_i$，给定环形序列可唯一构造闭游走。因此数环形配置等价于数闭游走。

\subsection{朴素思路：$26^{2026}$的枚举与为什么需要矩阵}

最直接的暴力算法是枚举全部$26^{2026}$种亮度分配，逐一验证三元组约束与$(2,6)$奇偶性，再用循环移位去重。$26^{2026}\approx 10^{2866}$远超任何计算机能处理的规模——整个可观测宇宙的原子数也仅约$10^{80}$。即便采用回溯法逐位搜索并利用三元组约束在每一位剪枝，$2026$位的搜索深度仍完全无法承受。然而约束的局部性——三元组只涉及连续三个灯珠，$(2,6)$计数可沿途累积——暗示状态空间可压缩为基于“最近两个灯珠”的马尔可夫链，进而用矩阵乘法批量计算所有路径。这正是从指数级暴力到多项式级代数计算的关键跳板。

\subsection{核心洞察：邻接矩阵的迹即环形序列计数}

将图中允许的转移记为$676\times 676$邻接矩阵$M$：$M_{ij}=1$若状态$i$可一步转移到状态$j$，否则为$0$。由矩阵乘法的组合解释，$(M\times M)_{ij}=\sum_k M_{ik}M_{kj}$枚举了从$i$经中间状态$k$到$j$的所有两步路径，恰好对应两步游走的组合计数。归纳地，$M^d$的第$(i,j)$元等于从$i$到$j$的长度为$d$的游走总数。

因此$\mathrm{tr}(M^d)=\sum_i (M^d)_{ii}$给出了所有长度为$d$的闭游走总数，记为$A(d)$。以一个$d=3$的具体环验证：设三个灯珠亮度为$(2,5,8)$，闭游走为$(2,5)\to(5,8)\to(8,2)\to(2,5)$。第一步检验$2+5+8=15\neq 26$，第二步检验$5+8+2=15\neq 26$，第三步检验$8+2+5=15\neq 26$。三步恰好覆盖环上全部$3$组连续三元组。对于$d=1$的平凡情形，$A(1)=\mathrm{tr}(M)$仅统计所有形如$(x,x)$的自环，要求$x+x+x\neq 26$。$3x=26$无整数解，故全部$26$个自环均合法，$A(1)=26$。

\subsection{符号矩阵：以$-1$权值从代数上分离$(2,6)$奇偶性}

$(2,6)$出现次数的奇偶性约束通过一个精巧的代数技巧解决。构造带符号矩阵$M'$：若某条转移的目标状态恰好为$(2,6)$（即$y=2,z=6$），则边权取$-1$（模意义下置为$\mathrm{MOD}-1$），否则取$1$。考虑任一长度为$d$的闭游走：目标为$(2,6)$的转移每出现一次贡献一个$-1$因子，其余转移贡献$1$，因此沿游走的全部边权乘积精确等于$(-1)^{\#(2,6)}$。这正是生成函数的思想——用$-1$给$(2,6)$的出现加权，在求和层面实现了奇偶性的差分编码。

故$B(d)=\mathrm{tr}((M')^d)$等于偶数情形数$C_0(d)$减去奇数情形数$C_1(d)$，即$B(d)=C_0(d)-C_1(d)$。同时$A(d)=\mathrm{tr}(M^d)=C_0(d)+C_1(d)$给出了不考虑奇偶性的所有合法环形序列数。两式相加得$2C_0(d)=A(d)+B(d)$，故$C_0(d)=(A(d)+B(d))/2$。模$998244353$下的除法乘以$2$的模逆元$(998244353+1)/2=499122177$实现。

\subsection{周期约束归约：何时用$A(d)$、何时用$C_0(d)$}

现在将视线转到Burnside引理下的$h(d)$定义上。$h(d)$是周期整除$d$的、长度为$N$的合法配置总数。周期为$d$的一个$N$长配置由一段长为$d$的“基本模式”重复$r=N/d$次铺排而成。在此重复构造下，环上全部三元组约束等价于要求该基本模式自身构成一个长度为$d$的合法环形序列——因为所有$d$个本质不同的三元组在重复次数增加时并不产生新的约束形式，且首尾衔接处也由基本模式的闭合性自动保证。

$(2,6)$的出现情况则需分情况讨论。设基本模式中$(2,6)$出现$t$次（$0\le t\le d$），则完整环上共出现$r\times t$次。若$r$为偶数，$r\times t$恒为偶数，与$t$无关，此时$h(d)$直接取$A(d)$；若$r$为奇数，则要求$t$本身为偶数，故需用$C_0(d)$。这一推理的优雅之处在于，它将一个全局约束（$N=2026$长度上的奇偶性）约化为了基本模式长度$d$上的局部约束，从而可以与矩阵幂的结果无缝对接。

\subsection{Burnside求和：四项因子的具体展开}

$N=2026=2\times 1013$，正因子集为$\{1,2,1013,2026\}$。四项的重复次数$r=N/d$依次为$2026,1013,2,1$。

由前述规则判定每一项的$h(d)$。$r=2026$为偶数$\Rightarrow h(1)=A(1)=26$。$r=1013$为奇数$\Rightarrow h(2)=C_0(2)$。$r=2$为偶数$\Rightarrow h(1013)=A(1013)$。$r=1$为奇数$\Rightarrow h(2026)=C_0(2026)$。

欧拉函数值的计算：$\varphi(2026)=\varphi(2)\varphi(1013)=1\times 1012=1012$（$1013$为素数故$\varphi(1013)=1012$）；$\varphi(1013)=1012$；$\varphi(2)=1$；$\varphi(1)=1$。

代入Burnside公式，所有项除以$N=2026$并取模：
\[
\mathrm{ans}=\frac{1012\cdot 26+1012\cdot C_0(2)+1\cdot A(1013)+1\cdot C_0(2026)}{2026}\pmod{998244353}.
\]
式中的除法通过乘以$2026$在模$998244353$下的逆元完成。由于$998244353$是素数且$2026$与之互素，该逆元必然存在。

\subsection{矩阵快速幂：二进制分解与访存优化}

为计算$A(1013)=\mathrm{tr}(M^{1013})$和$C_0(2026)$等项，需高效计算高次矩阵幂。直接迭代$1013$次矩阵乘法，单次$676^3\approx 3.09\times 10^8$次操作，总计$\sim 3\times 10^{11}$，不可接受。采用二进制快速幂策略：预处理$M^{2^k}$（$k=0,\dots,9$即$M^1$至$M^{512}$）共$9$次矩阵乘法。$1013$的二进制为$1111110101_2=512+256+128+64+32+16+4+1$，从$M^1$出发依次乘上$M^{2^2}$至$M^{2^9}$的$7$个已置位方幂，共$7$次乘法得到$M^{1013}$。

$\mathrm{tr}(M^{2026})$利用恒等式$\mathrm{tr}(M^{2026})=\sum_{i,j}(M^{1013})_{ij}(M^{1013})_{ji}$直接计算——对$M^{1013}$做$i,j$的二重循环累加$M_{ij}\cdot M_{ji}$即可，无需完整矩阵乘法。对$M'$同理，但$2026=1024+512+256+128+64+32+8+2$需预处理至$M'^{1024}$（$k=0,\dots,10$共$11$个方幂、$10$次乘法），再从$M'^2$出发乘上$7$个方幂得到$M'^{2026}$。

稠密矩阵乘法的性能瓶颈在于访存。为提升缓存局部性，预先转置右乘矩阵$B$为$B^T$。$C_{ij}=\sum_k A_{ik}B^T_{jk}$在内层对$k$的循环中，$A_{ik}$和$B^T_{jk}$均为行主序连续访问，大幅减少cache miss。累加过程使用\mil{__int128}类型容纳中间和——最坏情况约$676\times (998244352)^2\approx 6.76\times 10^{20}$，远低于\mil{__int128}上限$1.7\times 10^{38}$，遍历完一行$k$后对每个$C_{ij}$取模一次。对$M$和$M'$各自约需$17$次稠密矩阵乘法，总计约$34$次，每次$3.09\times 10^8$次基本操作，总操作量约$1.05\times 10^{10}$量级，在现代CPU上数秒内完成。空间需存储约$10$个$676\times 676$矩阵及对应转置，约$10\times 676^2\times 4\times 2\approx 36$MB，在合理范围内。

\subsection{代码实现}

本题为填空题，实际提交只需输出预计算的答案。但为了完整展示求解过程，以下同时给出幕后求解代码与提交代码。

\medskip\noindent 求解代码：
\inputminted{c++}{../src/p02_circular_light_ring_solve.cpp}

\medskip\noindent 提交代码：
\inputminted{c++}{../src/p02_circular_light_ring.cpp}

\section{题03：多项式积分（P16785）}

\subsection{问题重述}

题目给出一行字符串表示一个关于 $x$ 的一元多项式，要求输出其不定积分（不含积分常数 $C$）。输入多项式各项按次数从高到低排列，首项为正时前面不带 $+$ 号，首项为负时带 $-$ 号，后续各项以 $+$ 或 $-$ 分隔且符号归属于该项。系数为 $\pm 1$ 时可省略数字 $1$（例如 $x^{\wedge}3$ 表示 $1x^3$，$-x$ 表示 $-1x$），次数为 $1$ 时写作 $cx$ 或 $x$，次数为 $0$ 时写作整数常数，次数 $\ge 2$ 时写作 $cx^{\wedge}k$ 或 $x^{\wedge}k$。输出也需遵守类似格式，且系数为分数 $\frac{p}{q}$ 时要输出为 $p/q$，要求 $p,q$ 互质、$q>0$，负号放在分子上；若化简后分母为 $1$ 则直接输出整数。约束：项数 $\le 10^4$，次数 $\le 10^8$，系数绝对值 $\le 10^8$。

以样例 $1$ 为例，输入 \mil{4x\^{}5-4x\^{}3+2}，不定积分为 $\int(4x^5-4x^3+2)\,dx = \frac{4}{6}x^6 + \frac{-4}{4}x^4 + \frac{2}{1}x^1 = \frac{2}{3}x^6 - x^4 + 2x$，输出 \mil{2/3x\^{}6-x\^{}4+2x}。这个例子已经展示了问题的全部要素：系数从整数变成既约分数、分母为 $1$ 时省略分母、系数为 $-1$ 时省略数字 $1$。

\subsection{积分的线性性：逐项独立处理的数学依据}

不定积分的幂法则告诉我们 $\int cx^k \,dx = \frac{c}{k+1}x^{k+1} + C$。不定积分是一个线性算子，即 $\int (f+g) = \int f + \int g$ 且 $\int \alpha f = \alpha \int f$，因此多项式积分可以对每一项独立计算后直接拼接结果，各项之间互不干扰。更重要的是，每一项的次数加 $1$ 后，原来的降序排列自然保持——$k_1 > k_2$ 蕴含 $k_1+1 > k_2+1$，因此不需要额外排序。

于是整个问题分解为三个独立的阶段：将输入字符串解析为若干 $(c_i, k_i)$ 对；对每一对计算 $p_i = c_i,\; q_i = k_i+1,\; d_i = k_i+1$ 并用 $\gcd(|p_i|, q_i)$ 约简；按输出格式规则将结果拼接为字符串。数学核心只占三行代码，真正需要细致处理的是解析和格式化的诸多边界约定。

\subsection{解析的暗礁：隐式系数、符号归属与三种次数表示}

解析是本题最容易出错的部分。输入格式有若干“隐式约定”，需要逐字符扫描时精确判断。

符号归属规则：首项若为正则前面没有符号字符，首项若为负则有 $-$ 号；非首项前面一定有 $+$ 或 $-$。因此解析循环的每一轮，若当前字符是 $+$ 或 $-$ 则记录符号并前进指针，然后读取该项的系数和次数。

系数的隐式约定：若 $x$ 前面没有出现数字，则系数默认为 $1$（乘以符号后可能是 $-1$）。例如输入 \mil{x-1}，第一项在 \mil{x} 之前没有数字，系数为 $1$；第二项 \mil{-1} 有数字 $1$，系数为 $-1$。解析时用一个标志 \mil{have_coeff} 记录是否读到数字：读到了就用读到的值，没读到就默认为 $1$。读数字时逐位累加 \mil{coeff = coeff * 10 + digit}。

次数的三种表示：若当前字符是 \mil{x}，说明该项含有变量；接着若下一个字符是 \mil{\^{}} 则后面跟有指数数字（次数 $\ge 2$），否则次数为 $1$；若根本没遇到 \mil{x}，说明该项是常数，次数为 $0$。例如输入 \mil{4x\^{}5-4x\^{}3+2}，第一项的 \mil{x} 后跟 \mil{\^{}5}，次数为 $5$；最后一项 \mil{+2} 没有 \mil{x}，次数为 $0$。而 \mil{x} 单独出现时（如样例 $4$ 的 \mil{x-1}），第一项遇到 \mil{x} 且后面没有 \mil{\^{}}，次数为 $1$。

解析过程的正确性可以用几个边界输入检验：常数多项式 \mil{7} 应解析为 $(7,0)$；纯 $x$ 项 \mil{x} 应解析为 $(1,1)$；负系数无数字 \mil{-x} 应解析为 $(-1,1)$；大系数大次数如 \mil{-100x\^{}100} 应解析为 $(-100,100)$。解析结果的 \mil{vector<pair<ll,ll>>} 按输入顺序保存，即已经按次数降序排列。

\subsection{分数化简：符号归一化与分母恒正的约定}

对每一项 $(c_i, k_i)$，积分后的分子为 $c_i$，分母为 $k_i+1$，次数为 $k_i+1$。由于输入保证 $k_i \ge 0$，分母至少为 $1$。

约分步骤：设 $g = \gcd(|c_i|, k_i+1)$，用欧几里得算法在 $\mathrm{O}(\log \min(|c_i|, k_i+1))$ 时间内求出，然后分子分母同除以 $g$。注意 $\gcd$ 的参数应取绝对值，因为 $c_i$ 可能为负。

符号归一化步骤：输出规则要求分母恒为正，负号放在分子上。约分后若分母为负，则将分子分母同时取反。实现中只需判断 \mil{if (den < 0) \{ num = -num; den = -den; \}}。

以样例 $2$ 的第一项 $-100x^{100}$ 为例：$c = -100$，$k = 100$，分母 $101$，$\gcd(100, 101) = 1$（$100$ 和 $101$ 互质），约分后分子 $-100$，分母 $101$，次数 $101$。第二项 $105x^2$：$c=105$，$k=2$，分母 $3$，$\gcd(105,3)=3$，约分后分子 $35$，分母 $1$，次数 $3$——分母为 $1$ 表示系数化简为整数。两项组合输出为 \mil{-100/101x\^{}101+35x\^{}3}。

\subsection{输出格式的考验：首项符号、系数省略与分数表示}

输出阶段需要逐项判断三项内容：该项前面的符号、系数部分的字符串、变量部分的字符串。

符号规则分情况：若为第一项，正系数前面不加 $+$，负系数前面加 $-$；若为非首项，正系数加 $+$，负系数加 $-$。实现中只需检查 \mil{num} 的正负即可。

系数部分需要决定是否省略数字 $1$。规则是“单项系数为 $\pm 1$ 时省略数字 $1$”——这仅当分母为 $1$ 且分子绝对值为 $1$ 时才触发。例如积分结果中可能出现 $1x^3$，应输出为 \mil{x\^{}3}；$-1x$ 应输出为 \mil{-x}。对于分数系数（分母 $>1$），即使分子绝对值为 $1$，也不能省略，必须完整输出分子和分母。因为 $1/2$ 不是 $\pm 1$，省略 $1$ 会变成 \mil{/2}，无意义。因此代码逻辑为：若 \mil{den == 1} 且 \mil{abs_num == 1}，省略数字；若 \mil{den == 1} 且 \mil{abs_num != 1}，输出 \mil{abs_num}；若 \mil{den > 1}，输出 \mil{abs_num/den}。

变量部分：次数为 $1$ 时仅输出 \mil{x}（不输出指数 $^{\wedge}1$）；次数 $\neq 1$ 时输出 \mil{x\^{}deg}。注意常数项（次数 $0$）经过积分后次数变为 $1$，因此积分结果中不存在次数为 $0$ 的项——常数项的积分变成了 $x$ 的一次项，这正是样例 $3$ 中 $7$ 积分得到 $7x$ 的原因。

\subsection{复杂度分析与实现细节}

设输入字符串长度为 $n$，项数为 $m$（$m \le 10^4$）。每个字符最多被扫描一次，解析阶段 $\mathrm{O}(n)$。积分阶段每项做一次 $\gcd$，欧几里得算法复杂度 $\mathrm{O}(\log \max(|c_i|, k_i))$，由于 $\max(|c_i|, k_i) \le 10^8$，单次 $\gcd$ 约 $27$ 次迭代，总计 $m \times \log 10^8 \approx 2.7 \times 10^5$ 次运算。输出阶段 $\mathrm{O}(m)$。总体时间复杂度 $\mathrm{O}(n + m \log M)$，在给定约束下完全可行。

空间复杂度 $\mathrm{O}(m)$，存储解析结果和积分结果两组 \mil{vector}，每组约 $10^4$ 个元素，每个元素含三个 \mil{long long}，总计不足 $1$ MB。

实现中几个细节值得注意。使用 \mil{long long}（$64$ 位有符号整数）而非 \mil{int}，因为系数可达 $10^8$、次数可达 $10^8$，在计算 \mil{den = deg + 1} 和 \mil{gcd(num, den)} 时不会溢出，但若涉及乘法（如输出前做格式拼接）时 $64$ 位也足够。$\gcd$ 函数应先对参数取绝对值，因为分子可能为负，而欧几里得算法在负数上的行为是实现定义的。解析函数中标志 \mil{have_coeff} 至关重要：它区分了“系数为 $1$（省略数字）”和“常数的系数恰好是 $1$”这两种情形——对于输入 \mil{x}，没有数字但系数应为 $1$；对于输入 \mil{1}，有数字且系数就是 $1$。两种情形解析结果一致，但逻辑路径不同。最后，输出循环中注意首项与非首项的符号处理差异，以及分子绝对值 \mil{abs_num} 在符号已输出的前提下应使用绝对值。

\subsection{代码实现}
\inputminted{c++}{../src/p03_polynomial_integral.cpp}

\section{题04：安全路径（P16786）}

\subsection{问题建模：子树奇偶性与有序点对计数}

题目给出一棵以 $1$ 为根的 $n$ 个节点的无向树。
记 $\mathit{sz}[x]$ 为以 $x$ 为根的子树节点总数（含 $x$ 自身）。
若 $\mathit{sz}[x]$ 为偶数，称 $x$ 为平稳节点。
对于不同节点 $x,y$，若树中唯一简单路径上每个节点都平稳，
则有序对 $(x,y)$ 是一条安全路径，$(x,y)$ 与 $(y,x)$ 视为两条不同路径。
输入 $n$ 和 $n-1$ 条边，输出安全路径总数，$n \le 500000$，
要求线性或近线性时间算法。
本质是在树上按局部性质筛选节点子集，再统计子集内满足全平稳条件的点对数量。

\subsection{暴力枚举的局限：$\mathrm{O}(n^2)$ 在 $n=5\times10^5$ 下的困境}

最朴素思路：$\mathrm{O}(n)$ 算出 $\mathit{sz}$ 和平稳性，
再枚举 $\mathrm{O}(n^2)$ 个有序点对，逐对验证路径节点平稳性。
即便用 LCA 优化路径提取，单次验证最坏扫 $\mathrm{O}(n)$ 个节点，
总复杂度 $\mathrm{O}(n^3)$。即使能做到 $\mathrm{O}(1)$ 回答任意点对是否安全，
$n^2 \approx 2.5\times10^{11}$ 次查询本身已严重超时。
$n=5\times10^5$ 下，必须放弃显式枚举，寻找结构性计数方法。

\subsection{核心洞察：平稳节点诱导子图是森林}

重新审视条件：若 $x$ 到 $y$ 路径全平稳，则路径上相邻节点对也必平稳。
整条路径完全位于平稳节点诱导子图中——仅保留平稳节点，
且只保留两端均平稳的边。树连通无环，诱导子图必为森林。

在此森林中，$x$ 与 $y$ 之间有安全路径当且仅当同属一个连通分量：
原树两点简单路径唯一，全平稳则留在诱导子图中，两点连通；
连通路径也必全平稳。因此，大小为 $s$ 的连通分量内，
不同两点均构成两条有序安全路径，贡献 $s\cdot(s-1)$。
答案是所有分量贡献之和。

以示例验证：$n=6$，边 $(1,2),(1,3),(1,6),(2,4),(3,5)$，根 $1$。
$\mathit{sz}[4]=\mathit{sz}[5]=\mathit{sz}[6]=1$（奇，非平稳），
$\mathit{sz}[2]=\mathit{sz}[4]+1=2$（偶，平稳），
$\mathit{sz}[3]=\mathit{sz}[5]+1=2$（偶，平稳），
$\mathit{sz}[1]=\mathit{sz}[2]+\mathit{sz}[3]+\mathit{sz}[6]+1=6$（偶，平稳）。
平稳集 $\{1,2,3\}$。两端均平稳的边 $(1,2),(1,3)$ 形成大小 $3$ 的分量，
贡献 $3\times2=6$，与示例吻合。

\subsection{算法推导：自底向上DP中的分量合并与屏障结算}

可在诱导子图上 DFS 得分量大小，但显式构造需额外筛边。
更优雅的做法是原树上自底向上 DP，边算分量边累加答案，
全程仅对树两次遍历。

第一步：BFS 得层级与子树大小。从根 BFS，得 $\mathit{parent}[u]$ 和层次序
$\mathit{order}$（根到叶）。逆序处理：将 $\mathit{sz}[u]$ 累至 $\mathit{sz}[\mathit{parent}[u]]$，
完成后 $\mathit{stable}[u]=(\mathit{sz}[u]\bmod2=0)$。

第二步：自底向上 DP 算分量。再次逆序处理节点 $u$，
令 $\mathit{comp}[u]$ 为 $u$ 子树内以 $u$ 为顶的平稳连通分量大小。
若 $u$ 非平稳，$u$ 成屏障——向祖先的路径被 $u$ 中断。
遍历子 $v$：若 $v$ 平稳，分量无法再上延，结算
$\mathit{ans}\gets\mathit{ans}+\mathit{comp}[v]\cdot(\mathit{comp}[v]-1)$。
若 $u$ 平稳，$\mathit{cur}=1$，遍历子 $v$：
$v$ 平稳则边 $(u,v)$ 两端均平稳，合并 $\mathit{cur}\gets\mathit{cur}+\mathit{comp}[v]$；
$v$ 非平稳则其分量已在 $v$ 处结算，跳过。$\mathit{comp}[u]=\mathit{cur}$。

全节点处理完后，若根 $1$ 平稳，其分量未被子孙屏障结算，
手动追加 $\mathit{ans}\gets\mathit{ans}+\mathit{comp}[1]\cdot(\mathit{comp}[1]-1)$。

以链状树演示：$n=4$，边 $(1,2),(2,3),(3,4)$，根 $1$。
$\mathit{sz}[4]=1$（奇），$\mathit{sz}[3]=2$（偶），$\mathit{sz}[2]=3$（奇），$\mathit{sz}[1]=4$（偶）。
平稳 $\{1,3\}$，被非平稳 $2$ 隔开。
BFS 序 $[1,2,3,4]$，逆序：
$4$ 叶且非平稳，无操作；
$3$ 平稳，$\mathit{cur}=1$，子 $4$ 非平稳，$\mathit{comp}[3]=1$；
$2$ 屏障，子 $3$ 平稳，$\mathit{ans}+=1\times0=0$；
$1$ 平稳，$\mathit{cur}=1$，子 $2$ 非平稳，$\mathit{comp}[1]=1$；
结算根 $\mathit{ans}+=1\times0=0$。答案为 $0$，符合预期。

\subsection{复杂度与实现：线性扫描、long long 与迭代遍历}

两次遍历：BFS 建序 $\mathrm{O}(n)$，逆序累加 $\mathit{sz}$ 和 DP 各 $\mathrm{O}(n)$，
每条边各访问一次，总时间 $\mathrm{O}(n)$。空间 $\mathrm{O}(n)$（邻接表加五个线性数组），
$n=500000$、256 MB 下绰绰有余。

实现要点。答案须用 64 位有符号整数（C++ \mil{long long}）：
最坏全树平稳且连通，$n(n-1)\approx2.5\times10^{11}$ 远超 32 位 \mil{int} 范围。
树深可达 $n$（退化链），递归 DFS 会栈溢出，必用迭代 BFS。
$\mathit{order}$ 逆序即自底向上序。邻接表 \mil{vector<vector<int>>}，
读边双向建边，BFS 用 $\mathit{parent}$ 防回访。
结算公式 $s(s-1)$ 以 64 位运算：\mil{1LL * comp[v] * (comp[v] - 1)}。

\subsection{代码实现}
\inputminted{c++}{../src/p04_safe_paths.cpp}

\section{题05：生态廊道（P16787）}

\subsection{问题建模：路径图上局部重排的边权最大化}

题目给出了$n$个排成一列的监测位点，每个位点$i$拥有一个基因多样性编码$a_i$（$0\le a_i\le 10^9$）。定义交互强度为所有相邻位点编码按位异或之和：$S=\sum_{i=1}^{n-1}(a_i\oplus a_{i+1})$。我们可以将这条序列建模为一条含有$n$个顶点的路径图，第$i$条边连接顶点$i$和$i+1$，其权值即为$a_i\oplus a_{i+1}$，总权值$S$便是交互强度。

工程师的优化操作是：在全序列上选择两个互不重叠、长度均为$3$的连续区域（只能选择一次），对每个区域内部的三个位点进行任意重新排列。排列只改变数值在区域内的出现顺序，不改变数值集合本身。记两个区域的起点下标分别为$i$和$j$（$1\le i,j\le n-2$），约束$|i-j|\ge 3$保证了两个长度为$3$的区域不会交叉或重叠。我们的目标是在所有满足条件的$(i,j)$以及所有可能的排列中，最大化最终的交互强度。

为了与代码保持一致，下文统一采用0-indexed下标（数组下标范围为$[0,n-1]$）。此时区间起点$i$（$0\le i\le n-3$）所覆盖的三个位置为$\{i,i+1,i+2\}$。可选的合法起点总数为$m=n-2$。原始交互强度$S_{\text{base}}=\sum_{k=0}^{n-2}(a_k\oplus a_{k+1})$在输入后即可$\mathrm{O}(n)$求出，所有优化方案的总强度均可写为$S_{\text{base}}+\Delta$的形式，其中$\Delta$为排列操作带来的总增量。于是问题等价于求$\Delta$的最大值。

\subsection{朴素暴力：从三次方到局部增量的第一步收缩}

最直接的暴力解法分三层：第一层枚举区间起点对$(i,j)$，共有$\mathrm{O}(n^2)$种有效组合；第二层，每个区间有$3!=6$种排列，组合起来共$36$种方案；第三层，对每种方案重新计算全部$n-1$条边的异或和，耗时$\mathrm{O}(n)$。总复杂度$\mathrm{O}(36n^3)$，当$n=2\times10^5$时$n^3\approx8\times10^{15}$，完全不可行。

稍加观察即可发现优化空间：一个长度为$3$的区间被重新排列后，受影响的边总共不超过四条——区间内部有两条边$(i,i+1)$和$(i+1,i+2)$，区间左侧最多有一条跨越边界的边$(i-1,i)$（仅当$i>0$时存在），区间右侧最多也有一条跨越边界的边$(i+2,i+3)$（仅当$i+2<n-1$即$i<n-3$时存在）。序列中其余所有边的贡献均不发生变化。因此对每组$(i,j,p,q)$，我们可以仅重算这至多$8$条受影响边的贡献变化，在$\mathrm{O}(1)$时间内完成增量的计算，整体复杂度降为$\mathrm{O}(36n^2)$。然而$n=2\times10^5$时$n^2\approx4\times10^{10}$，依然远超时限。问题的关键在于：我们不应该先确定$(i,j)$再计算增量，而应当反过来——先独立计算出每个区间各自的最优增量，再利用结构性质高效地合并它们。

\subsection{核心洞察：区间增量可独立预计算，间距决定合并方式}

考察单个区间：设0-indexed下标$i$为起点，原值$(a_i,a_{i+1},a_{i+2})=(x,y,z)$。该区间的某种排列$p$将三个位置上的值分别变为$b_0,b_1,b_2$（三者是$(x,y,z)$的一个置换）。该排列对交互强度的增量$\Delta_i(p)$由三部分组成：
\[
\begin{aligned}
\Delta_i^{\text{left}}(p) &=
\begin{cases}
(a_{i-1}\oplus b_0)-(a_{i-1}\oplus x), & i>0,\\
0, & i=0;
\end{cases}\\[4pt]
\Delta_i^{\text{in}}(p) &= (b_0\oplus b_1)+(b_1\oplus b_2)-(x\oplus y)-(y\oplus z);\\[4pt]
\Delta_i^{\text{right}}(p) &=
\begin{cases}
(b_2\oplus a_{i+3})-(z\oplus a_{i+3}), & i<n-3,\\
0, & i=n-3.
\end{cases}
\end{aligned}
\]
三者之和即为区间$i$在排列$p$下的总增量$\Delta_i(p)$。对每个起点$i$，枚举全部$6$种排列即可在$\mathrm{O}(1)$时间内得到该处的最优独立增量$\Delta_i^{\max}=\max_p\Delta_i(p)$。

这些“独立增量”的隐含假设是：区间两端的邻居位点均保持原始值不变。当两个所选区间距离足够远（中间至少隔了一个未被调整的位置），这一假设对两个区间同时成立——它们各自的受影响边集合互不相交，总增量可以直接相加：$\Delta=\Delta_i(p)+\Delta_j(q)$。然而，若两个区间恰好紧邻（即$j=i+3$），则区间$i$的右边界和区间$j$的左边界恰好共享同一条边$(i+2,i+3)$，两套独立增量都试图修改这条边、却各自假设对端为原值，从而导致叠加出错。这一观察自然将问题划分为两种情形：独立情形（$|i-j|\ge 4$）和共享边界情形（$|i-j|=3$）。

\subsection{独立情形：前缀后缀最大值将配对从O(n²)压缩至O(n)}

当$|i-j|\ge 4$时，不妨设$j\ge i+4$。由于区间$i$覆盖$[i,i+2]$，区间$j$覆盖$[j,j+2]$，位置$i+3$位于两者之间且未被任何区间修改——它恰好是区间$i$右边界所使用的“原始邻居”，也是区间$j$左边界所依赖的那个位置的左侧元素。两个区间的受影响边集合完全不相交，故总最优增量为$\Delta_i^{\max}+\Delta_j^{\max}$。问题转化为：在所有满足$|i-j|\ge 4$的起点对中，求$\Delta_i^{\max}+\Delta_j^{\max}$的最大值。

这一子问题通过前缀最大值和后缀最大值在$\mathrm{O}(m)$时间内解决。定义前缀数组$\mathit{pref}[k]=\max_{0\le t\le k}\Delta_t^{\max}$和后缀数组$\mathit{suff}[k]=\max_{k\le t<m}\Delta_t^{\max}$，均为$\mathrm{O}(m)$建表。对每个起点$i$，所有满足$j\le i-4$的$j$中最优者即为$\mathit{pref}[i-4]$（当$i\ge4$时有效），所有满足$j\ge i+4$的$j$中最优者即为$\mathit{suff}[i+4]$（当$i+4<m$时有效）。遍历全部$i$，检查$\Delta_i^{\max}+\mathit{pref}[i-4]$和$\Delta_i^{\max}+\mathit{suff}[i+4]$，取全局最大值即覆盖了所有独立情形的最优配对。

以样例验证：$n=7$，$a=[1,2,3,4,5,6,7]$（1-indexed），则$m=5$个0-indexed起点的原始编码分别为：$i=0$覆盖$[1,2,3]$，$i=1$覆盖$[2,3,4]$，$i=2$覆盖$[3,4,5]$，$i=3$覆盖$[4,5,6]$，$i=4$覆盖$[5,6,7]$。对于$i=1$，排列$[2,4,3]$给出内部增量$(2\oplus4)+(4\oplus3)-(2\oplus3)-(3\oplus4)=6+7-1-7=5$，右边界增量$(3\oplus5)-(4\oplus5)=6-1=5$，故$\Delta_1^{\max}=10$；$i=4$的最优排列即原序，$\Delta_4^{\max}=0$。注意$|1-4|=3$，恰为共享边界情形，独立情形的前缀/后缀扫描不会覆盖它。

\subsection{共享边界：精确修正重叠边的重复调整}

当$|i-j|=3$时，由对称性只需考虑$j=i+3$（另一方向在枚举中自动覆盖）。区间$i$覆盖$[i,i+1,i+2]$，区间$j$覆盖$[i+3,i+4,i+5]$，两者仅通过边$(i+2,i+3)$相连。记区间$i$在排列$p$下的右端值为$R_i(p)$（即原来的$b_2$），区间$j$在排列$q$下的左端值为$L_j(q)$（即$c_0$），原始共享边的贡献$\mathit{orig}=a_{i+2}\oplus a_{i+3}$。直接相加$\Delta_i(p)+\Delta_j(q)$会将同一原始贡献扣减两次、又将各自基于“对端不变”假设的新贡献各加一次，形成偏差。正确的合并增量需要先卸掉两区间各自对共享边的独立修正，再补上两个新区间端点直接异或的真实贡献：
\[
\Delta_{\text{combined}}(p,q)=\Delta_i(p)+\Delta_j(q)-\Delta_i^{\text{right}}(p)-\Delta_j^{\text{left}}(q)+(R_i(p)\oplus L_j(q))-\mathit{orig}.
\]

回到样例验证此公式。取$i=1$、$j=4$，$R_1(p)=3$（排列$[2,4,3]$的右端值），$L_4(q)=5$（原始排列的左端值），$\mathit{orig}=a_3\oplus a_4=4\oplus5=1$。已知$\Delta_1(p)=10$、$\Delta_1^{\text{right}}(p)=5$，$\Delta_4(q)=0$、$\Delta_4^{\text{left}}(q)=0$。代入得$10+0-5-0+(3\oplus5)-1=5+6-1=10$。加上$S_{\text{base}}=16$后总强度为$26$，与样例答案一致。

对于每一对满足$j=i+3$的起点，枚举两个区间各自的$6$种排列（共$36$种组合），用上述公式计算合并增量并维护全局最大值。值得注意的是，这里不能简单地取$\Delta_i^{\max}+\Delta_j^{\max}$再作边界修正——因为使$\Delta_i$最大的那个排列$p$，其右端值$R_i(p)$不一定与使$\Delta_j$最大的那个排列$q$的左端值$L_j(q)$搭配时得到最大的$(R_i\oplus L_j)$，所以必须遍历全部$36$种组合。好在这一枚举每对只需常数时间，总开销$\mathrm{O}(36n)$完全可接受。

\subsection{整体复杂度与实现要点}

算法分为三个阶段，均为线性。预处理阶段：对$m=n-2$个起点各枚举$6$种排列，每个排列在$\mathrm{O}(1)$时间内完成三条增量的计算，总耗时$\mathrm{O}(6m)=\mathrm{O}(n)$。独立情形阶段：构建$\mathit{pref}$和$\mathit{suff}$数组各需一次$\mathrm{O}(m)$扫描，配对扫描又为$\mathrm{O}(m)$。共享边界阶段：枚举$\mathrm{O}(n)$对紧邻区间，每对枚举$36$种排列组合，总耗时$\mathrm{O}(36n)$。因此总体时间复杂度$\mathrm{O}(n)$，空间复杂度$\mathrm{O}(n)$（主要为存储每起点的$6$组增量分量和端点值）。对于$n=2\times10^5$的数据量，约$7.2\times10^6$次共享边界枚举运算，在现代CPU上毫秒级可完成。

实现时需注意几个细节。异或值的累加可能达到$2\times10^5\cdot 2^{30}\approx 2\times10^{14}$，超出的\mil{int}的$2^{31}-1\approx2.1\times10^9$，必须使用64位有符号整数（C++中的\mil{long long}）。预处理时需为每个起点的每种排列同时保存总增量\mil{total_delta}、左边界增量\mil{delta_left}、右边界增量\mil{delta_right}以及左右端点实际值\mil{left_val}和\mil{right_val}，供共享边界阶段取用。边界位置的增量为零需显式判定：起点$i=0$时左边界分量记为$0$，$i=m-1$（即区间右端已触碰序列末尾）时右边界分量记为$0$。共享边界情形中，将$a_{i+2}$与$a_{i+3}$的异或值作为$\mathit{orig}$直接代入公式即可，无需引入额外的索引映射。

\subsection{代码实现}
\inputminted{c++}{../src/p05_ecology_corridor.cpp}

\section{题06：独立三角形（P16788）}

\subsection{问题重述}

二维平面上有 $N$ 个互不重合的整数点，每个点的纵坐标只可能是 $1$ 或 $-1$。我们需要用这些点连出尽可能多的三角形，每个三角形须满足三个顶点不共线，且不同三角形之间不能共享顶点——一个点一旦作为三角形的顶点，就已经连接到该三角形内部的另外两个顶点，共消耗了两条线段的端点，这也意味着每个点至多参与一个三角形（因为每个点最多作为两条线段的端点）。在三角形数量取到最大的所有方案中，进一步要求所有三角形的面积之和最大。输入保证该最大面积和为整数。$N \le 2 \times 10^5$，$|x_i| \le 10^9$。

\subsection{建模：当所有点只分布在两条水平线上}

这是一个典型的“受限组合优化”问题。输入本质上是两组一维坐标的集合：上方点集（$y = 1$）的横坐标列表和下方点集（$y = -1$）的横坐标列表。设上方点个数为 $A$，下方点个数为 $B$，则 $A + B = N$。

由于所有点只能位于 $y = 1$ 和 $y = -1$ 两条直线上，而三个顶点若共线则不能构成三角形，因此任何合法三角形的三个顶点不可能全在上方或全在下方。每个三角形的顶点构成只可能有两种情形：两个上方点搭配一个下方点（以下称为“以底在上”的三角形），或一个上方点搭配两个下方点（“以底在下”的三角形）。两种情形中，三角形的面积公式都是 $\frac12 \times \text{底边长度} \times 2 = \text{底边长度}$，因为底边由同侧的两个顶点确定，其对侧顶点到底边的垂直距离恒为 $2$。换言之，三角形的面积恰好等于其底边两端点的横向距离，与第三条顶点的具体位置无关。

\subsection{朴素思路：暴力枚举为何不可行}

如果采用最暴力的做法，可以从 $N$ 个点中枚举所有可能的三元组，检查其是否构成合法三角形，再在合法三角形集合中找出若干个两两不共享顶点的方案。枚举所有三角形的时间至少是 $\mathrm{O}(N^3)$，在 $N = 2 \times 10^5$ 的规模下完全不可行。即便我们利用“所有点仅在两条线上”这一性质把候选三角形压缩到“二上一下”和“一上二下”两类，每类的候选数量依然是 $\mathrm{O}(N^2)$ 级别——上方点的每对组合都可以搭配任意一个下方点。因此必须寻找更深刻的结构性结论，将指数级的组合选择转化为简单的线性规划与贪心。

\subsection{核心洞察：每个三角形消耗一上一下两条边}

设我们最终构造了 $x$ 个以底在上的三角形和 $y$ 个以底在下的三角形，总三角形数 $T = x + y$。从顶点消耗的角度看：每个以底在上的三角形消耗 $2$ 个上方点和 $1$ 个下方点，每个以底在下的三角形消耗 $1$ 个上方点和 $2$ 个下方点。因此上方点的总消耗为 $2x + y$，下方点的总消耗为 $x + 2y$，它们分别不能超过 $A$ 和 $B$：

\begin{equation}\label{eq:constraints}
2x + y \le A, \qquad x + 2y \le B, \qquad x \ge 0,\; y \ge 0.
\end{equation}

将两式相加得到 $3(x + y) \le A + B = N$，即 $T = x + y \le \lfloor N/3 \rfloor$。另一方面，每个三角形无论哪种类型都至少需要 $1$ 个上方点和 $1$ 个下方点，所以 $T \le \min(A, B)$。综合这两条独立的上界，我们得到三角形数量的理论上界：

\begin{equation}\label{eq:maxT}
T_{\max} = \min\!\big(\lfloor N/3 \rfloor,\; A,\; B\big).
\end{equation}

这个上界有没有意义？以样例 $1$ 为例，$N = 3$，三个点分别为 $(0, 1), (2, 1), (1, -1)$，则 $A = 2$，$B = 1$。$T_{\max} = \min(1, 2, 1) = 1$，确实可以达到 $1$ 个三角形。样例 $2$ 中 $N = 6$，上方点坐标 $0, 1, 3$（$A = 3$），下方点坐标 $1, 2, 4$（$B = 3$），$T_{\max} = \min(2, 3, 3) = 2$，同样可达到。

\subsection{上界的可达性证明：线性不等式组的整数解}

上界给出的是必要条件，我们还需证明它总是充分的——即对于任意满足 $T = T_{\max}$ 的 $T$，总存在非负整数 $x, y$ 使得 $x + y = T$ 且满足不等式组~(\ref{eq:constraints})。将 $y = T - x$ 代入，约束转换为 $x$ 的一个闭区间：

\[
2x + (T - x) \le A \;\Longrightarrow\; x \le A - T,
\]
\[
x + 2(T - x) \le B \;\Longrightarrow\; x \ge 2T - B.
\]

再加上 $0 \le x \le T$，得到 $x$ 的可行区间：

\begin{equation}\label{eq:xrange}
\max(0,\, 2T - B) \le x \le \min(T,\, A - T).
\end{equation}

该区间非空等价于左右端点满足 $\max(0,\, 2T - B) \le \min(T,\, A - T)$。分情况验证：若 $T = \lfloor N/3 \rfloor$ 且 $T \le A, B$，由 $3T \le A + B$ 可推出 $2T - B \le A - T$，同时 $A - T \ge 0$（因为 $T \le A$），且 $2T - B \le T$（因为 $T \le B$），故区间非空。若 $T = A$（此时 $A$ 是三者中最小者），代入得 $T = A \le \lfloor N/3 \rfloor$，结合 $A + B = N$ 可知 $2A \le B$，此时可行区间退化为 $\max(0, 2A - B) \le x \le \min(A, 0)$，即 $x = 0$——全部三角形均为以底在下的类型。对称地，若 $T = B$ 则 $x = T$——全部为以底在上的类型。由此可见，$T_{\max}$ 总是可以通过合理分配两种三角形的数量来实现。

\subsection{最大化面积之和：底边配对的最优策略}

由于每个三角形的面积等于其底边两端点的横向距离，而第三条顶点的横坐标不影响面积，因此最大化面积之和的问题可以拆解为：从 $A$ 个上方点中选出 $2x$ 个点，两两配成 $x$ 条底边，使这 $x$ 条底边长度之和最大；同时对下方点做同样操作（选出 $2y$ 个点配成 $y$ 条底边）。剩下的 $A - 2x$ 个上方点和 $B - 2y$ 个下方点充当第三条顶点，它们的横坐标无关紧要。

现在问题简化为：给定 $m$ 个已排序的点（横坐标从小到大），要从中选出 $2k$ 个点配成 $k$ 对，使得每对的横向距离之和最大。直觉上，要使距离最大，应该让选出的点尽可能分散——最左和最右配对、次左和次右配对，依此类推。这一直觉可以用重排不等式严格证明：对于已排序数组 $p_0 \le p_1 \le \dots \le p_{m-1}$，在所有选出 $2k$ 个元素并配对的方案中，配对方式 $(p_0, p_{m-1}), (p_1, p_{m-2}), \dots, (p_{k-1}, p_{m-k})$ 使得 $\sum (p_{m-i} - p_i)$ 取到最大值。原因在于：内层的点无论如何选，其贡献的距离不会超过外层点产生的跨度；如果我们从两端选取最远的 $k$ 对，中间的点即使被选为第三条顶点，也不会影响底边长度，而任何试图使用中间点作为底边端点的方案都会使距离和变小。

以样例 $2$ 的上方点集 $(0, 1, 3)$ 为例：$A = 3$，若 $x = 1$，需选出 $2$ 个点配对。选取最左的 $0$ 和最右的 $3$，底边长度为 $3$；若选 $0, 1$ 或 $1, 3$，长度仅为 $1$ 或 $2$，均不如 $3$。

基于这一结论，我们分别将上方点和下方点按横坐标升序排列后，预计算“前缀最大宽度和”：

\[
\mathit{topW}[k] = \sum_{i=0}^{k-1} (t_{A-1-i} - t_i),\qquad
\mathit{botW}[k] = \sum_{i=0}^{k-1} (b_{B-1-i} - b_i).
\]

其中 $\mathit{topW}[k]$ 表示从上方点中选出 $2k$ 个并以最优方式配成 $k$ 条底边所能获得的最大底边长度之和，$\mathit{botW}[k]$ 含义对称。注意 $\mathit{topW}[k]$ 仅在 $2k \le A$ 时有定义（同理 $\mathit{botW}[k]$ 要求 $2k \le B$）；超出此范围的 $k$ 取值为 $0$（无法构造那么多以该侧为底的三角形）。

有了这两个前缀和数组，对于一个可行的 $(x, y)$ 组合（其中 $y = T - x$，$x$ 满足区间~(\ref{eq:xrange})），该方案的总面积即为 $\mathit{topW}[x] + \mathit{botW}[y]$。遍历 $x$ 的所有可行取值并记录最大值，即可得到在 $T = T_{\max}$ 个三角形的前提下所能获得的最大面积之和。

\subsection{复杂度分析}

整个算法的瓶颈在于排序：分别对上方点集和下方点集按横坐标排序，时间复杂度为 $\mathrm{O}(A \log A + B \log B) = \mathrm{O}(N \log N)$。前缀和数组的计算是线性的 $\mathrm{O}(\max T)$，遍历 $x$ 的可行区间同样是 $\mathrm{O}(\max T) \le \mathrm{O}(N)$。在 $N \le 2 \times 10^5$ 的数据范围下，$\mathrm{O}(N \log N)$ 的排序足以在时限内完成。空间复杂度为 $\mathrm{O}(N)$，用于存储两个坐标数组和两个前缀和数组。

\subsection{实现细节}

代码实现中有几个值得注意的细节。首先，由于 $|x_i| \le 10^9$ 且 $N$ 可达 $2 \times 10^5$，面积之和可能超过 $32$ 位整型的表示范围，因此前缀和数组 $\mathit{topW}$ 和 $\mathit{botW}$ 以及最终答案 $\mathit{bestArea}$ 必须使用 \mil{long long}（$64$ 位有符号整数）。其次，在计算前缀和时，需要确保索引不越界：\mil{topW[k]} 仅当 $2k \le A$ 且 $k \le \max T$ 时才可计算；代码中通过 \mil{min(A/2, maxT)} 来控制循环上界，对于超出此范围的 $k$，$\mathit{topW}[k]$ 和 $\mathit{botW}[k]$ 保持初始化时的 $0$，这恰好对应“无法构造 $k$ 个以该侧为底的三角形则贡献为 $0$”的语义。此外，$T_{\max}$ 可能为 $0$（例如所有点都在同一条线上），此时需特判并输出 \mil{0 0}，避免后续计算中出现空的循环范围导致未定义行为。最后，$x$ 的可行区间由 $T_{\max}$ 和 $A, B$ 共同确定：下界 $\max(0, 2T - B)$ 保证下方点够用，上界 $\min(T, A - T)$ 保证上方点够用；这两条边界均来源于不等式组~(\ref{eq:constraints})的直接推导。

\subsection{代码实现}
\inputminted{c++}{../src/p06_independent_triangles.cpp}

\section{题07：核心任务调度（P16789）}

\subsection{问题重述}

有 $N$ 个任务，每个任务恰好需要一天完成，小蓝每天至多完成一个任务，也可以在某天不安排任何任务。天数从第 $1$ 天开始按正整数编号。每个任务 $i$ 有一个截止时间 $t_i$ 和一个价值 $w_i$：如果在第 $t_i$ 天或之前完成该任务，则获得价值 $w_i$；若超过截止时间，任务失效且不可再完成。此外，部分任务被标记为核心任务（$\text{is\_key}_i = 1$），其余为普通任务（$\text{is\_key}_i = 0$）。公司要求小蓝制定调度方案时，必须先保证完成的核心任务数量尽可能多，在所有满足这一要求的方案中再使总价值尽可能大。求最多能完成的核心任务数 $K$ 和该前提下的最大总价值 $V$。约束：$N \le 2 \times 10^5$，$t_i \le 10^9$，$w_i \le 10^9$。

样例中共 $4$ 个任务：任务 $1$ 的 $t=3, w=10$ 且为核心任务，任务 $2$ 的 $t=1, w=5$ 且为核心任务，任务 $3$ 的 $t=2, w=100$ 且为普通任务，任务 $4$ 的 $t=1, w=100$ 且为普通任务。最优方案为第 $1$ 天完成任务 $2$、第 $2$ 天完成任务 $3$、第 $3$ 天完成任务 $1$，完成 $2$ 个核心任务，总价值 $5+100+10=115$。

\subsection{从蛮力枚举到带权调度贪心}

先考虑最朴素的思路：枚举所有任务的子集，检查该子集是否存在一个可行调度（即能否将每个选中任务分配到一个不超过其截止时间且互不冲突的日期），若可行则记录其核心任务数和总价值，最终在所有可行子集中取字典序最优者。可行性检查本身就是一个经典问题：按截止时间排序后贪心分配最早可用日期，若所有选中任务都能分配到日期则可行。枚举 $2^N$ 个子集、每个子集花费 $\mathrm{O}(N \log N)$ 检查，总复杂度 $\mathrm{O}(2^N N \log N)$，在 $N = 2 \times 10^5$ 面前毫无希望。

但这道题的底层结构正是经典的带截止时间单位任务调度问题：$N$ 个任务各耗时一单位，各有截止时间和价值，要求选出一个可调度子集使总价值最大。该问题有一个优雅的拟阵贪心算法。将所有任务按截止时间从小到大排序，依次处理每个任务：将其价值加入一个小根堆，若此时堆的大小（即已选任务数）超过了当前任务的截止时间，说明时间槽不够用了，则弹出堆中价值最小的任务。处理完所有任务后，堆中剩余的即为所选任务。

以题目样例演示该算法（暂时忽略核心/普通的区分，仅按原始价值 $w_i$ 运行）。四个任务按截止时间排序后为：$(t=1,w=5)$、$(t=1,w=100)$、$(t=2,w=100)$、$(t=3,w=10)$。依次处理：第一个任务入堆，堆大小 $1 \le t=1$，保留；第二个任务入堆，堆大小 $2 > t=1$，弹出最小值 $5$，堆中只剩价值 $100$ 的任务；第三个任务入堆，堆大小 $2 \le t=2$，保留；第四个任务入堆，堆大小 $3 \le t=3$，保留。最终堆中有三个任务（价值 $100,100,10$），总价值 $210$。然而正确的最优答案是完成两个核心任务（价值 $5$ 和 $10$）加上普通任务 $3$（价值 $100$），总价值 $115$。可见直接最大化总价值的贪心无法保证核心任务数量优先。

\subsection{字典序双目标的数值编码}

本题的真正目标是一个字典序优化：第一关键字是核心任务数量 $K$（最大化），第二关键字是总价值 $V$（最大化）。等价地，我们希望找到一个方案使得二元组 $(K, V)$ 在字典序下最大。这提示我们可以将双目标编码为单一数值，从而复用经典的单目标贪心框架。

关键观察：设所有任务的原始价值之和为 $S = \sum_{i=1}^{N} w_i$，令 $C = S + 1$。定义每个任务的调整价值为
\[
\tilde{w}_i = w_i + C \cdot \text{is\_key}_i
\]
那么，对任意两个方案 $A$ 和 $B$，$\sum_{i \in A} \tilde{w}_i > \sum_{i \in B} \tilde{w}_i$ 当且仅当 $A$ 的字典序 $(K_A, V_A)$ 严格优于 $B$ 的字典序 $(K_B, V_B)$。原因如下：任何多完成一个核心任务的方案，其调整价值之和至少多出 $C$（因为新完成的核心任务贡献了 $C + w_i \ge C$）；而任何不改变核心任务数量的方案，其调整价值之和的差异至多为 $\sum w_i = S < C$。因此调整价值之和的比较天然地实现了先比较 $K$、再比较 $V$ 的字典序语义。

在样例中，$S = 10 + 5 + 100 + 100 = 215$，$C = 216$。四个任务的调整价值分别为 $\tilde{w}_1 = 10 + 216 = 226$，$\tilde{w}_2 = 5 + 216 = 221$，$\tilde{w}_3 = 100 + 0 = 100$，$\tilde{w}_4 = 100 + 0 = 100$。现在以调整价值 $\tilde{w}_i$ 替代原始价值 $w_i$ 运行前述带权调度贪心算法。

\subsection{贪心算法的逐步推演}

将四个任务按截止时间排序（截止时间相同时任意顺序均可，因为堆会在超出截止时间时自动淘汰价值最低者）：依次为 $(t=1, \tilde{w}=221)$、$(t=1, \tilde{w}=100)$、$(t=2, \tilde{w}=100)$、$(t=3, \tilde{w}=226)$。

第一步：处理截止时间 $1$、调整价值 $221$ 的任务 $2$。将其入堆，堆大小 $1 \le 1$，保留。此时堆中为 $\{221\}$。

第二步：处理截止时间 $1$、调整价值 $100$ 的任务 $4$。将其入堆，堆大小变为 $2$，超过了当前截止时间 $1$，触发淘汰。弹出堆中最小调整价值 $100$，堆中恢复为 $\{221\}$。

第三步：处理截止时间 $2$、调整价值 $100$ 的任务 $3$。将其入堆，堆大小 $2 \le 2$，保留。堆中为 $\{221, 100\}$。

第四步：处理截止时间 $3$、调整价值 $226$ 的任务 $1$。将其入堆，堆大小 $3 \le 3$，保留。堆中最终为 $\{221, 100, 226\}$，对应任务 $2$、$3$、$1$。

现在从最终堆中还原答案。核心任务数 $K$ 等于最终堆中核心任务的出现次数，在此为 $2$（任务 $2$ 和任务 $1$）。总价值 $V$ 等于最终堆中所有任务的原始价值之和，在此为 $5 + 100 + 10 = 115$。这与题目给出的最优解完全一致。

\subsection{为什么编码技巧是正确的}

上述编码技巧的正确性依赖于 $C > \sum w_i$ 这一不等式。直观上，$C$ 充当了“核心任务”这一维度的权重，使得任何核心任务数量的差异都能压倒总价值维度的任何差异。形式化地，考虑任意两个可行方案 $A$ 和 $B$，它们的调整价值总和之差为
\[
\Delta = \sum_{i \in A} \tilde{w}_i - \sum_{i \in B} \tilde{w}_i = C \cdot (K_A - K_B) + (V_A - V_B).
\]
若 $K_A > K_B$，则 $K_A - K_B \ge 1$，故 $\Delta \ge C - (V_B - V_A) \ge C - S = 1 > 0$，方案 $A$ 的调整价值更大。若 $K_A = K_B$，则 $\Delta = V_A - V_B$，调整价值的比较退化为纯价值比较。因此调整价值之和的最大化恰好等价于字典序 $(K, V)$ 的最大化。

值得注意的是，这一步编码并不改变贪心算法本身的正确性——带权调度贪心的正确性根植于该问题所对应的拟阵结构（独立集为所有存在可行调度的任务子集），调整价值仅是改变了每个元素的权重，拟阵的贪心算法对任意非负权重均成立。

\subsection{复杂度与实现细节}

整个算法分为三个阶段：计算权重和并构造调整价值，按截止时间排序，以及单次扫描加堆操作。第一阶段遍历所有任务一次，$\mathrm{O}(N)$；排序阶段 $\mathrm{O}(N \log N)$；扫描阶段每个任务至多入堆一次、出堆一次，堆操作 $\mathrm{O}(\log N)$，故 $\mathrm{O}(N \log N)$。总时间复杂度 $\mathrm{O}(N \log N)$，空间复杂度 $\mathrm{O}(N)$。在 $N = 2 \times 10^5$ 的约束下，$\log_2 N \approx 18$，总操作量约 $3.6 \times 10^6$ 级别，远在时限之内。

实现时有几个值得注意的细节。由于 $w_i$ 可达 $10^9$ 且 $N$ 可达 $2 \times 10^5$，$S = \sum w_i$ 可达 $2 \times 10^{14}$，需使用 \mil{int64_t}（即 \mil{long long}）存储 $C$ 和调整价值，避免溢出。排序时直接对截止时间排序即可，因为贪心正确性不依赖同截止时间内任务的处理顺序——堆会在总任务数超出截止时间时自动淘汰最劣者，无论它们以何种顺序入堆。另外，在扫描过程中，与其从最终的调整价值反推核心任务数和总价值，不如在入堆和出堆时直接维护 \mil{total_val} 和 \mil{core_cnt} 两个变量——入堆时加上对应任务的 $w_i$ 和 $\text{is\_key}_i$，出堆时减去——这样最终答案可以直接输出，无需额外解码步骤。堆中存储的是（调整价值，原始下标）的 pair，因为出堆时需要知道被淘汰的是哪个任务，以便正确更新 \mil{total_val} 和 \mil{core_cnt}。

\subsection{代码实现}
\inputminted{c++}{../src/p07_task_scheduler.cpp}

\section{题08：魔法前缀（P16790）}

\subsection{问题重述}

给定 $N$ 个小写字母单词，定义有效前缀集合为所有单词的全部非空前缀的集合（相同前缀去重）。小蓝和小桥轮流操作，小蓝先手。每次操作时，当前玩家从集合中选择一个非空前缀 $P$，然后删除集合中所有以 $P$ 为前缀的字符串（包括 $P$ 本身）。轮到某玩家时若集合为空则该玩家判负。两人均采取最优策略。每组测试用例判断先手小蓝（XiaoLan）还是后手小桥（XiaoQiao）必胜。约束：$T \le 10$，$N \le 10^5$，所有测试用例字符串总长度不超过 $10^6$。

\subsection{模型转化：Trie树上的公正组合游戏}

第一步是将字符串集合转化为Trie树。将所有单词插入一个Trie树，根结点（编号0）代表空前缀，不出现在游戏中；每个非根结点唯一对应一个有效前缀。于是游戏操作变为：选择任意一个非根结点，删除该结点及其整棵子树。被删除的子树从游戏中移除，但剩余部分可能分裂成多个互不影响的独立子树，每个子树以原根结点的某个未被删除的子结点为根。

例如，单词 \texttt{cat} 和 \texttt{car} 的Trie树包含：根结点（空），根的子结点 \texttt{c}，\texttt{c} 的子结点 \texttt{a}，\texttt{a} 的子结点 \texttt{t} 和 \texttt{r}。若某玩家选择结点 \texttt{a}，则 \texttt{a}、\texttt{t}、\texttt{r} 均被删除，仅剩 \texttt{c} 挂在根下。若该玩家选择的是 \texttt{c}，则整棵树被清空，游戏即刻结束。

这是一个公正组合游戏：双方可进行的操作完全相同，信息完全公开，胜负仅取决于状态而非玩家身份。根据Sprague-Grundy定理，每个状态对应一个Grundy数，终态（无任何结点的空集）Grundy为0；非终态的Grundy数等于其所有后继状态Grundy数的mex（最小非负未出现整数）。整个游戏由根结点的各个子结点为根的子树作为独立子游戏构成，总Grundy数等于各子游戏Grundy数的异或（Nim加法）和。若总Grundy数非零则先手必胜，否则后手必胜。

\subsection{朴素思路：穷举博弈树为何不可行}

最直接的想法是将每个有效前缀集合视为一个博弈状态，枚举当前玩家可选择的所有前缀，模拟双方最优走法。但状态数量随前缀集合大小呈指数增长：假设Trie有 $M$ 个非根结点，每次操作删除一个子树后剩余部分仍可能有多种组合方式，状态空间远远大于 $2^M$。而本题中总结点数可达 $10^6$ 级别，穷举博弈树完全不可行。

\subsection{子树Grundy递归公式的核心洞察}

关键突破在于观察操作的局部性质。在Trie树上，选择某个结点 $v$ 仅影响 $v$ 的子树内部结构，$v$ 的兄弟子树完全不受影响。换言之，对 $v$ 子树的操作是一个独立的子游戏，其Grundy数仅取决于 $v$ 下方树的结构，与整棵Trie的其余部分无关。

考虑以某个非根结点 $v$ 为根的子树。设 $v$ 有 $k$ 个子结点 $c_1, c_2, \dots, c_k$，以它们为根的子树的Grundy数分别为 $g_1, g_2, \dots, g_k$。我们想要求出 $v$ 的子树整体对应的Grundy数 $G(v)$。

当前玩家面对 $v$ 这棵子树时有两种选择：一是直接选择 $v$ 本身，将整棵子树清空，后继状态Grundy为 $0$。二是在某个子结点 $c_i$ 的子树内部进行操作，子游戏 $c_i$ 的Grundy数从 $g_i$ 变为其某个后继值 $h$（由Grundy数的定义，$c_i$ 的所有后继Grundy数恰好覆盖 $0, 1, \dots, g_i - 1$ 中至少这些值），而其他子游戏保持不变。此时整体Grundy变为 $1 + h \oplus \bigoplus_{j \ne i} g_j$。这里加 $1$ 是因为对 $c_i$ 子树的操作发生于 $v$ 的子树内部，$v$ 本身仍存在于游戏中；从博弈代数角度看，$v$ 与 $k$ 个子游戏之间是“串联”关系：$v$ 的一次额外选择权叠加在 $k$ 个并列子游戏之上。

可以严格证明 $G(v) = 1 + \bigoplus_{i=1}^k g_i$。下面用一个三结点小例来展示直观含义：考虑一条长度为 $2$ 的链 \texttt{root} -- \texttt{a} -- \texttt{b}。叶结点 \texttt{b} 无子结点，$G(\texttt{b}) = 1 + 0 = 1$。结点 \texttt{a} 有一个子结点 \texttt{b}（其 $g = 1$），故 $G(\texttt{a}) = 1 + 1 = 2$。这意味着：若玩家直接选 \texttt{a}，游戏结束（Grundy 0）；若玩家选 \texttt{b}，则 \texttt{a} 仍在场，剩下一个Grundy为 $1$ 的子树，后继Grundy为 $1$。可达集合为 $\{0, 1\}$，mex 正好是 $2$。公式成立。

\subsection{算法推导：自底向上计算总Grundy值}

算法流程如下。首先建Trie：将所有单词插入一棵Trie树，每个结点维护26个子结点指针。由于Trie的建树顺序保证子结点编号必然大于父结点，我们可以按结点编号从大到小遍历（自底向上），对每个非根结点 $v$：
\[
G(v) = 1 + \bigoplus_{\substack{c \in \{0,\dots,25\} \\ \text{child}(v, c) \,\neq\, \text{null}}} G(\text{child}(v, c))
\]
叶结点无子结点，异或和为 $0$，故 $G(\text{leaf}) = 1$。

以样例为例验证公式。第一组样例：单词 \texttt{ab} 与 \texttt{cd}。Trie根有两个子结点 \texttt{a} 和 \texttt{c}，各为一条长度 $2$ 的链。叶结点 \texttt{b} 的 $G = 1$，\texttt{a} 的 $G = 1 + 1 = 2$；同理 \texttt{d} 的 $G = 1$，\texttt{c} 的 $G = 2$。总Grundy $= 2 \oplus 2 = 0$，后手必胜，输出 XiaoQiao，与样例一致。

第二组样例：单词 \texttt{cat} 与 \texttt{car}。叶结点 \texttt{t} 的 $G = 1$，叶结点 \texttt{r} 的 $G = 1$。结点 \texttt{a} 有两个子结点，$G(\texttt{a}) = 1 + (1 \oplus 1) = 1 + 0 = 1$。结点 \texttt{c} 有子结点 \texttt{a}（$G = 1$），$G(\texttt{c}) = 1 + 1 = 2$。根下仅一个子结点 \texttt{c}，总Grundy $= 2 \ne 0$，先手必胜，输出 XiaoLan，也与样例一致。从博弈角度看，小蓝第一步选择 \texttt{c} 即可清空整棵树，小桥无子可走。

\subsection{复杂度分析}

建Trie和自底向上计算Grundy值均仅对每个结点访问常数次。Trie的总结点数为所有单词的字符总数，即不超过 $10^6$。每个结点的Grundy计算只需遍历26个子结点并做异或，总操作量 $\mathrm{O}(|\Sigma| \cdot M)$，其中 $|\Sigma| = 26$，$M$ 为总结点数。整体时间复杂度 $\mathrm{O}(\sum |s_i|)$，空间复杂度同为 $\mathrm{O}(\sum |s_i|)$。对总字符数 $10^6$ 的输入，所有操作在毫秒级即可完成。

\subsection{实现细节}

在代码实现中，Trie使用扁平数组存储而非动态分配：\mil{nxt[node * 26 + c]} 表示从当前结点沿字符 $c$ 走到的子结点编号，\mil{-1} 表示无边。该设计避免了大量指针分配，访问速度更快，也使得结点按自然插入顺序获得递增编号，天然满足子结点编号大于父结点的性质。

自底向上遍历时，从最大编号结点递减至根结点（编号1）。根结点（编号0）对应空前缀，不参与Grundy值计算；其各子结点的Grundy值异或即为总Grundy值。总Grundy非零时先手小蓝必胜，否则后手小桥必胜。

\subsection{代码实现}
\inputminted{c++}{../src/p08_magic_prefix.cpp}

\section{题09：神秘排列（P16791）}

\subsection{从输入到数学结构：问题建模}

输入是一个正整数 $n \le 10^7$，要求统计 $1$ 到 $n$ 的所有 $n!$ 个排列中有多少个是“好排列”。一个排列 $a_1, a_2, \ldots, a_n$ 中，位置 $i$ 称为平衡位置，当且仅当左侧比 $a_i$ 小的元素个数 $L_i = |\{j < i : a_j < a_i\}|$ 恰好等于右侧比 $a_i$ 大的元素个数 $R_i = |\{j > i : a_j > a_i\}|$。如果一个排列的平衡位置数量不少于 $\lceil n/2 \rceil$，则称该排列为好排列。以 $n = 4$ 为例，$\lceil 4/2 \rceil = 2$，也就是说至少需要 $2$ 个平衡位置才算是好排列。排列 $(4, 3, 2, 1)$ 中，位置 $1$ 的左侧没有元素故 $L_1 = 0$，且 $a_1 = 4$ 是最大值，右侧没有比 $4$ 大的元素故 $R_1 = 0$，位置 $1$ 平衡；类似可验证其余三个位置也都平衡，因此该排列是好排列。本题的核心挑战在于 $n$ 可以高达 $10^7$，枚举排列完全不可能，必须找到平衡位置的结构性特征，将问题转化为可计算的组合计数。

\subsection{暴力枚举为何行不通}

最直接的朴素思路是：枚举所有 $n!$ 个排列，对每个排列扫描每个位置计算 $L_i$ 和 $R_i$，统计平衡位置数量，判断是否为好排列。对于 $n = 4$ 来说，$4! = 24$ 个排列尚可手动检验，但 $n = 10$ 时 $10! \approx 3.6 \times 10^6$ 已经接近暴力极限，而 $n = 5000$ 时 $n!$ 是一个天文数字。显然需要将“平衡位置”这一直观定义转化为等价的代数条件，把计数问题从排列层面转化为组合公式层面。

\subsection{平衡位置的充要条件：$a_i = n - i + 1$}

我们从单个位置 $i$ 出发推导平衡位置的等价条件。设位置 $i$ 处的值为 $a_i$。在整个 $1 \ldots n$ 的排列中，严格小于 $a_i$ 的数共有 $a_i - 1$ 个，严格大于 $a_i$ 的数共有 $n - a_i$ 个（没有等于 $a_i$ 的数，因为排列中元素互异）。这些数分布在位置 $i$ 的左右两侧：小于 $a_i$ 的数中有 $L_i$ 个在左侧，剩余 $(a_i - 1) - L_i$ 个在右侧；大于 $a_i$ 的数中有 $R_i$ 个在右侧，剩余 $(n - a_i) - R_i$ 个在左侧。位置 $i$ 左侧恰好有 $i - 1$ 个位置，每个位置要么放了一个小于 $a_i$ 的数，要么放了一个大于 $a_i$ 的数，因此
\[
L_i + \bigl(n - a_i - R_i\bigr) = i - 1.
\]
如果位置 $i$ 是平衡位置，则 $L_i = R_i$，代入上式消去 $L_i$ 和 $R_i$：
\[
L_i + n - a_i - L_i = i - 1 \quad\Longrightarrow\quad n - a_i = i - 1 \quad\Longrightarrow\quad a_i = n - i + 1.
\]
这就是平衡位置的必要条件：若 $i$ 平衡，则 $a_i$ 必须恰好等于 $n - i + 1$。例如 $n = 4$ 时，若位置 $1$ 平衡，必须 $a_1 = 4$；若位置 $2$ 平衡，必须 $a_2 = 3$；若位置 $3$ 平衡，必须 $a_3 = 2$；若位置 $4$ 平衡，必须 $a_4 = 1$。

反过来，假设 $a_i = n - i + 1$ 成立。此时 $a_i$ 是确定的数值，我们来验证 $L_i = R_i$ 必然成立。由 $a_i = n - i + 1$ 可知小于 $a_i$ 的数共有 $n - i$ 个，大于 $a_i$ 的数共有 $i - 1$ 个。设左侧有 $x$ 个小于 $a_i$ 的数，即 $L_i = x$。左侧共有 $i - 1$ 个位置，其中除了 $x$ 个小于 $a_i$ 的数之外，其余 $(i - 1) - x$ 个位置填的是大于 $a_i$ 的数。而全部 $i - 1$ 个大于 $a_i$ 的数必须分配在左右两侧——左侧有 $(i - 1) - x$ 个，右侧自然有 $R_i = (i - 1) - ((i - 1) - x) = x$ 个。于是 $R_i = x = L_i$，充分性得证。因此，位置 $i$ 是平衡位置当且仅当 $a_i$ 恰好等于 $n - i + 1$。这个条件非常简洁：平衡位置就是那些“值与位置和为 $n + 1$”的位置。

\subsection{变换视角：不动点计数问题}

上述条件 $a_i = n - i + 1$ 启发我们引入变换 $q_i = n + 1 - a_i$。由于 $a$ 是 $[1, n]$ 的排列，$q$ 也是 $[1, n]$ 的一个排列，且 $a \leftrightarrow q$ 是一一映射。此时，
\[
a_i = n - i + 1 \iff n + 1 - a_i = i \iff q_i = i.
\]
也就是说，原排列 $a$ 中的一个平衡位置恰好对应变换后排列 $q$ 中的一个不动点——即 $q_i = i$，元素值等于其下标的位置。原问题随之转化为：在 $n$ 个元素的所有排列中，不动点个数不少于 $\lceil n/2 \rceil$ 的排列有多少个？

这是一个经典的组合计数问题。设 $f(n, m)$ 表示 $n$ 元排列中恰有 $m$ 个不动点的排列个数。要得到恰好 $m$ 个不动点，先从 $n$ 个位置中选出哪 $m$ 个位置是不动点，有 $\binom{n}{m}$ 种选择方式；选定后这 $m$ 个位置的值就已经确定（$q_i = i$），而剩余的 $n - m$ 个位置上不能出现任何不动点，即这 $n - m$ 个位置上的子排列必须是一个错位排列，记 $D(k)$ 为 $k$ 元错位排列数（即没有任何元素在其原始位置的排列数）。于是
\[
f(n, m) = \binom{n}{m} \cdot D(n - m).
\]
好排列的总数即为所有满足 $m \ge \lceil n/2 \rceil$ 的 $f(n, m)$ 之和：
\[
\mathrm{ans} = \sum_{m = \lceil n/2 \rceil}^{n} \binom{n}{m} D(n - m).
\]

\subsection{错位排列的通项：容斥原理推导}

$D(k)$ 的标准公式可以通过容斥原理推导。$k$ 个元素的全排列有 $k!$ 个，从中减去至少有一个不动点的排列，加回至少有两个不动点的排列，以此类推。固定某 $i$ 个位置为不动点后，其余 $k - i$ 个位置可任意排列，有 $(k - i)!$ 种方式，而选择哪 $i$ 个固定位置有 $\binom{k}{i}$ 种，于是
\[
D(k) = \sum_{i=0}^{k} (-1)^i \binom{k}{i} (k - i)!
     = k! \sum_{i=0}^{k} \frac{(-1)^i}{i!}.
\]
例如 $D(0) = 1$（空排列），$D(1) = 0$（单个元素必然在原始位置），$D(2) = 1$（交换两个元素），$D(3) = 2$（三轮换及其反向），$D(4) = 9$，与直观一致。

\subsection{化简求和：$\mathrm{O}(n)$ 的线性计算}

直接按 $\mathrm{ans} = \sum_{m=k}^{n} \binom{n}{m} D(n - m)$ 计算，每个 $m$ 需要 $\mathrm{O}(n - m)$ 时间求 $D(n - m)$ 的和式，总复杂度为 $\mathrm{O}(n^2)$，在 $n = 10^7$ 时不可接受。我们将 $\binom{n}{m} D(n - m)$ 展开并化简：
\begin{align*}
\binom{n}{m} D(n - m)
    &= \frac{n!}{m!\,(n - m)!} \cdot (n - m)! \sum_{i=0}^{n - m} \frac{(-1)^i}{i!} \\
    &= \frac{n!}{m!} \sum_{i=0}^{n - m} \frac{(-1)^i}{i!} \\
    &= n! \cdot \frac{1}{m!} \cdot F(n - m),
\end{align*}
其中 $F(t) = \sum_{i=0}^{t} (-1)^i / i!$（在模 $10^9 + 7$ 意义下，除法用乘法逆元实现）。令 $k = \lceil n/2 \rceil$，则
\[
\mathrm{ans} = n! \sum_{m=k}^{n} \frac{1}{m!} \cdot F(n - m) \pmod{10^9 + 7}.
\]
这个形式可以仅用一趟扫描完成计算。首先用单变量累积计算 $n! \bmod M$；然后利用费马小定理 $\mathit{inv\_fact}[n] = (n!)^{M - 2} \bmod M$ 求阶乘逆元，再倒推 $\mathit{inv\_fact}[i - 1] = \mathit{inv\_fact}[i] \cdot i \bmod M$ 得到所有阶乘逆元（无需存储完整的阶乘数组，因为中间的 $\mathit{fact}[t]$ 可通过 $\mathit{inv\_fact}[t]$ 和已知的 $n!$ 按需获得，而此处求和式中恰不需要它们）。接着构建前缀和数组 $F$：$F[0] = 1$，对 $t = 1, 2, \ldots, n - k$，$F[t] = F[t - 1] + (-1)^t \cdot \mathit{inv\_fact}[t] \pmod M$，其中奇数 $t$ 时 $(-1)^t \equiv M - 1$。最后遍历 $m = k, k + 1, \ldots, n$，累加 $\mathit{fact\_n} \cdot \mathit{inv\_fact}[m] \cdot F[n - m] \pmod M$ 即得答案。

以 $n = 4$ 为例验证：$k = 2$，$n! = 24$。$\mathit{inv\_fact}[0] = 1$，$\mathit{inv\_fact}[1] = 1$，$\mathit{inv\_fact}[2] = 500000004$，$\mathit{inv\_fact}[3] = 166666668$，$\mathit{inv\_fact}[4] = 41666667$。$F[0] = 1$，$F[1] = 1 + (M - 1) \cdot 1 \equiv 0$，$F[2] = 0 + 1 \cdot \mathit{inv\_fact}[2] = 500000004$。$m = 2$ 项：$24 \cdot 500000004 \cdot F[2] \bmod M$；$m = 3$ 项：$24 \cdot 166666668 \cdot F[1] = 0$；$m = 4$ 项：$24 \cdot 41666667 \cdot F[0] = 1$。三项之和为 $7$，与样例一致。

\subsection{复杂度分析}

预处理 $n!$ 需要 $\mathrm{O}(n)$ 时间；用快速幂求 $\mathit{inv\_fact}[n]$ 需 $\mathrm{O}(\log M)$ 时间（约 $30$ 次乘法）；倒推所有阶乘逆元需要 $\mathrm{O}(n)$ 时间；计算 $F$ 数组需要 $\mathrm{O}(n - k) = \mathrm{O}(n)$ 时间；最终累加也需要 $\mathrm{O}(n)$ 时间。总体时间复杂度为 $\mathrm{O}(n)$。在 $n = 10^7$ 时，约 $10^7$ 次模乘和模加，在现代 CPU 上可在数百毫秒内完成。空间方面，$\mathit{inv\_fact}$ 数组尺寸为 $n + 1 \approx 10^7$，以 \mil{int} 存储约占 $40$ MB；$F$ 数组尺寸为 $n - k + 1 \approx n/2 \approx 5 \times 10^6$，约占 $20$ MB；两项合计约 $60$ MB，在蓝桥杯国赛的通常内存限制（$256$ MB 或 $512$ MB）下完全可行。

\subsection{实现细节与注意事项}

代码实现中几个值得关注的要点。第一，模数 $M = 10^9 + 7$ 是一个质数，这保证了费马小定理求逆元的正确性。第二，$n$ 可能为 $1$（此时 $\lceil 1/2 \rceil = 1$，只有排列 $[1]$ 是好的，答案为 $1$），代码中的循环在 $n = 1$ 时 $k = 1$，$\mathit{fact\_n} = 1$，$F$ 数组长度为 $0$，求和仅含 $m = 1$ 一项，结果为 $1$，边缘情况自动处理正确。第三，$\mathit{inv\_fact}$ 数组需要用到 \mil{vector<int>}，空间占用在 C++ 中为 $n \cdot \texttt{sizeof(int)}$，对于 $n = 10^7$ 为 $40$ MB，注意不要额外存储 $\mathit{fact}$ 数组；$F$ 数组的长度仅为 $n - k + 1$，可以用较小的空间分配。第四，累加过程中模运算密集，使用 \mil{long long} 做中间乘法以避免溢出，这是 $10^9 + 7$ 级别模运算的常规做法。

\subsection{代码实现}

\inputminted{c++}{../src/p09_mysterious_permutation.cpp}

\section{题10：不互质游戏（P16792）}

\subsection{问题重述与博弈建模}

桌面上有 $n$ 个正整数 $a_1, \dots, a_n$，两人轮流操作。每次操作选择一个数 $x$，将其替换为更小的 $y$，要求 $y < x$ 且 $\gcd(x, y) > 1$。无法操作者输。多组数据，$T \le 200$，每组 $n \le 1000$，$a_i \le 300000$，所有数据的 $n$ 之和不超过 $100000$。

这是一个公平组合游戏：双方可执行的操作完全相同，不存在随机因素，局面信息完全公开；每个桌面数字构成一个独立的子游戏，整个游戏是这些子游戏的直和。根据 Sprague–Grundy 定理，直和局面的 Grundy 数等于各子游戏 Grundy 数的异或和，先手必胜当且仅当该异或和非零。因此核心任务是为每个可能的 $a_i$（即 $1$ 到 $300000$ 的所有整数）预计算其子游戏的 Grundy 数 $G(a_i)$，其中 $G(x)$ 定义为：
\[
G(x) = \operatorname{mex}\{\, G(y) \mid 1 \le y < x,\ \gcd(x, y) > 1 \,\}.
\]
这里 $\operatorname{mex}$（minimum excluded value）表示集合中未出现的最小非负整数。边界：$G(1) = 0$（$1$ 无合法后继），对任意质数 $p$ 也有 $G(p) = 0$（不存在小于 $p$ 且与 $p$ 不互质的数）。

\subsection{朴素思路——暴力枚举为何不可行}

最直接的递推思路是：对每个 $x = 2, 3, \dots, 300000$，枚举所有 $y < x$，检查 $\gcd(x, y) > 1$ 是否成立，收集符合条件的 $G(y)$ 再求 $\operatorname{mex}$。对于 $x = 300000$，需要检查约 $3 \times 10^5$ 个候选 $y$，累计约 $\frac{1}{2} \times (3 \times 10^5)^2 \approx 4.5 \times 10^{10}$ 次 $\gcd$ 调用。即使辗转相除法每次只需 $\mathrm{O}(\log M)$ 时间，总运算量已远超 $10^{11}$ 量级，在任何评测环境下都无法通过。

可以做初步优化：不枚举所有 $y$，而是枚举 $x$ 的每个奇质因子 $p$ 的倍数。设 $x$ 的相异奇质因子为 $p_1, \dots, p_k$，合法 $y$ 的集合恰为 $\bigcup_i \{\, p_i, 2p_i, 3p_i, \dots \mid p_i \cdot t < x \,\}$。但最坏情况下 $x$ 含有多个小质因子时合法 $y$ 的数量仍可达到 $\mathrm{O}(x)$。例如 $x = 2310 = 2 \times 3 \times 5 \times 7 \times 11$，仅质因子 $3$ 就贡献了近 $770$ 个合法后继。更重要的是，我们需要的并非 $y$ 本身，而是 $G(y)$ 的集合——这一观察暗示我们可以“按质因子汇总 Grundy 值”，这正是后续位集优化的核心动机。

\subsection{偶数的封闭公式——归纳与上界证否}

偶数情形最为规整。设 $x = 2k$。所有更小的偶数 $2, 4, \dots, 2(k-1)$ 均与 $x$ 共享因子 $2$，因此全部构成合法后继。递推前几项：$G(2) = \operatorname{mex}\varnothing = 0$；$G(4) = \operatorname{mex}\{0\} = 1$；$G(6) = \operatorname{mex}\{0, 1\} = 2$；$G(8) = \operatorname{mex}\{0, 1, 2\} = 3$。归纳可知前 $k-1$ 个偶数的 Grundy 数恰好为 $0, 1, \dots, k-2$，因此偶数后继覆盖了所有 Grundy 值从 $0$ 到 $k-2$ 的连续区间。

接下来须论证 Grundy 值 $k-1$ 是否可达。若存在某个奇数 $y < 2k$ 使得 $G(y) = k-1$，则 $G(2k)$ 至少为 $k$，公式将不成立。关键观察：对任意正整数 $y$，恒有 $G(y) \le y/2 - 1$——$\operatorname{mex}$ 不可能超过集合的大小，而 $y$ 的后继 Grundy 值集合的规模至多为 $y/2 - 1$。对任意不超过 $2k-1$ 的奇数 $y$，$G(y) \le (2k-1)/2 - 1 = k - 1.5$，取整为 $k-2$；对小于 $2k$ 的偶数 $y$，$G(y) = y/2 - 1 \le (2k-2)/2 - 1 = k-2$。因此 $k-1$ 永不可达。于是得到简洁的封闭公式：
\[
G(2k) = k - 1 = \frac{x}{2} - 1.
\]
验证 $x = 12 = 2 \times 6$：公式给出 $G(12) = 5$。偶数后继 $2, 4, 6, 8, 10$ 贡献 $\{0, 1, 2, 3, 4\}$；奇数后继 $3$（$\gcd=3$）和 $9$（$\gcd=3$）贡献 $\{0, 1\}$；全体集合为 $\{0, 1, 2, 3, 4\}$，$\operatorname{mex}=5$，吻合。

\subsection{含因子 $3$ 的奇数的公式——从小数枚举中发现模式}

奇数质数的 Grundy 数为 $0$ 是平凡的。对于奇数合成数，先考察含因子 $3$ 的情况。手工计算前几项：$G(3) = 0$（质数）。$x = 9 = 3^2$：合法后继为 $3, 6$，$G(3) = 0$，$G(6) = 2$，$\operatorname{mex}\{0, 2\} = 1$。$x = 15 = 3 \times 5$：合法后继有 $3, 5, 6, 9, 10, 12$，Grundy 值为 $\{0, 0, 2, 1, 4, 5\}$，$\operatorname{mex} = 3$。$x = 21 = 3 \times 7$：后继 $3, 6, 7, 9, 12, 14, 15, 18$ 对应 Grundy $\{0, 2, 0, 1, 5, 6, 3, 8\}$，$\operatorname{mex} = 4$。$x = 27 = 3^3$：后继涵盖 $3$ 和 $2$ 的所有 $3$ 的倍数达到基本全覆盖，经过细致枚举可得 $\operatorname{mex} = 6 = \lfloor 24/4 \rfloor$。

将这些结果排成一列：$G(9) = \lfloor 6/4 \rfloor = 1$，$G(15) = \lfloor 12/4 \rfloor = 3$，$G(21) = \lfloor 18/4 \rfloor = 4$，$G(27) = \lfloor 24/4 \rfloor = 6$。模式清晰浮现：
\[
G(x) = \left\lfloor \frac{x - 3}{4} \right\rfloor, \qquad \text{对奇数 $x$ 且 $3 \mid x$}.
\]
该公式的深层逻辑是：$3$ 的倍数的 Grundy 值分布形成了缺口恰在 $\lfloor (x-3)/4 \rfloor$ 处的连续前缀。即使 $x$ 额外含有其他质因子（如 $15$ 含 $5$，$21$ 含 $7$，$45$ 含 $5$），这些额外质因子的倍数所提供的 Grundy 值全部落在上述前缀内部，永远无法精准地填补缺口。其根本原因是一条关于互质性的不动点原理：所有 Grundy 值为 $\lfloor (x-3)/4 \rfloor$ 的数必然与 $3$ 互质（否则它本该由某个 $3$ 的倍数更早取得），而既然它与 $3$ 互质，任何含有因子 $3$ 的数 $x$ 都无法将它作为合法后继。

\subsection{一般奇合成数——位集维护与按质分治}

对于最小质因子至少为 $5$ 的奇合成数（如 $25, 35, 49, 55, 77, \dots$），没有统一的封闭公式，必须回到递推框架。令 $P = \{p_1, \dots, p_k\}$ 为 $x$ 的所有互异奇质因子。由于所有 $y < x$ 的 $G(y)$ 均已计算完毕，从 $x$ 可达的 Grundy 值集合恰为“每个 $p \in P$ 的全部已处理倍数所取得的 Grundy 值的并集”。于是我们为每个奇质数 $p$ 维护一个集合 $\mathit{reach}[p]$，处理 $x$ 时求所有 $\mathit{reach}[p]$（$p \in P$）的并集的 $\operatorname{mex}$，处理完再将 $G(x)$ 纳入各 $\mathit{reach}[p]$。

对所有 $25997$ 个奇质数各维护独立集合将导致巨额开销，必须按质数大小分治。将奇质数分为“小质数”（$\le 1000$）和“大质数”（$> 1000$）两类。$M = 300000$ 内约 $168$ 个小质数，每个的倍数多达 $M/p \ge 300$ 个，适合用 \mil{std::bitset<MAX_G>} 维护——每个 \mil{bitset} 约 $19$ KB，总内存约 $3$ MB。合并时直接按位或，利用 $64$ 位机器字并行，$k$ 个 \mil{bitset} 的合并代价约为 $\mathit{MAX\_G} \cdot k / 64$ 次字操作。大质数约有 $25600$ 个，但其倍数极度稀疏（$M/p \le 300$），对每个大质数以 \mil{std::vector<pair<int,int>>} 按处理顺序记录（倍数，Grundy 值），合并时顺序扫描并筛选 \mil{multiple < x} 的条目追加到局部位集。

以 $x = 25 = 5^2$ 为例演示递推过程。处理 $x = 25$ 时，$\operatorname{spf}(25) = 5$（小质数），奇质因子仅有 $\{5\}$。此时 $\mathit{reach}[5]$ 中已有 $5$ 的倍数 $5, 10, 15, 20$ 的 Grundy 值：$G(5)=0$，$G(10)=4$，$G(15)=3$，$G(20)=9$。因此局部位集中置位的比特为 $0, 3, 4, 9$，从 $0$ 开始搜索，遇到未置位的最小值 $1$，故 $G(25) = 1$。处理完 $25$ 后，将比特 $1$ 置入 $\mathit{reach}[5]$。再处理 $x = 35 = 5 \times 7$ 时，奇质因子为 $\{5, 7\}$（$7$ 也是小质数），局部位集取 $\mathit{reach}[5] \mid \mathit{reach}[7]$，搜索 $\operatorname{mex}$ 即得 $G(35)$。整个过程对每个 $x$ 的代价正比于其奇质因子个数乘上小质数 \mil{bitset} 的长度或大质数向量的已扫描条目数。

\subsection{预处理流程与复杂度分析}

完整预处理的执行流程如下。第一步，线性筛：对 $i = 2$ 到 $M$，若 $\operatorname{spf}[i] = 0$ 则标记其为质数并加入质数列表；用每个已发现的质数 $p$ 筛去 $i \cdot p$，同时赋予 $\operatorname{spf}$ 值，当 $p > \operatorname{spf}[i]$ 时跳出内层循环以保证线性。同时将每个奇质数分配为小质数索引或大质数索引。该步骤的复杂度为 $\mathrm{O}(M \log \log M)$，对 $M = 300000$ 耗时远小于 $0.01$ 秒。

第二步，递推计算 $G(x)$。对于偶数 $x$，$G(x) = x/2 - 1$；对于奇质数 $x$，$G(x) = 0$；对于含因子 $3$ 的奇合成数，$G(x) = \lfloor (x-3)/4 \rfloor$；这三种情况各自只需 $\mathrm{O}(1)$ 时间，覆盖了约 $85\%$ 的数。对于 $\operatorname{spf} \ge 5$ 的奇合成数（约 $2 \times 10^4$ 个），每组 $\omega(x)$ 个奇质因子（$\omega(x)$ 最多约 $6$，均摊远小）需要做 \mil{bitset} 或操作或向量扫描。小质数 \mil{bitset} 按位或的单词量约为 $\mathit{MAX\_G} / 64 \approx 2300$ 字；大质数向量的扫描总条目数不超过 $\sum_{p > 1000} M/p \approx M(\ln \ln M - \ln \ln 1000) \approx 6 \times 10^5$ 个。综合运算量在 $10^7$ 级，完全可在 $0.2$ 秒内完成。

第三步，对每组查询读入 $n$ 个数，计算 $\bigoplus_i G(a_i)$。若异或和非零则输出 \mil{"Yes"}，否则输出 \mil{"No"}。$\sum n \le 100000$ 保证所有查询的累加只需线性时间。

空间复杂度：小质数 \mil{bitset} 数组约 $3$ MB；\mil{spf} 数组和 $G$ 数组各约 $1.2$ MB（\mil{int} 型）；大质数向量总条目数不超过各组互异奇质因子的出现次数之和，约在百万量级，每个条目 $8$ 字节故约 $8$ MB；总计不超过 $15$ MB。

\subsection{实现细节与注意事项}

线性筛内层循环中，乘法 \mil{i * p} 需使用 \mil{long long} 比较以防止 \mil{int} 溢出（$300000 \times 300000 \approx 9 \times 10^{10}$ 超出了 $32$ 位有符号整数的范围）。提取 $x$ 的互异奇质因子时，反复除以 $\operatorname{spf}[tmp]$ 并在 \mil{last_p} 与当前质因子相同时跳过，这样既去除了因子 $2$ 又自动去重。大质数的 \mil{reachable_large} 向量中条目按倍数递增自然有序，这是因为倍数（即各 $x$）按处理序递增；扫描时检查 \mil{entry.multiple >= x} 直接 \mil{break}，避免了无效遍历。局部位集 \mil{seen} 在每轮循环中通过 \mil{seen.reset()} 清零，清空一个长 \mil{bitset} 的代价足以忽略。$\operatorname{mex}$ 由循环 \mil{while (seen[g]) g++;} 得到，由于绝大多数情况下 $\operatorname{mex}$ 值不大，循环迭代次数极短。最终的胜负判定须牢记 SG 定理的结论：异或和非零则先手必胜；若所有 $a_i$ 的 Grundy 值异或为零，则后手有必胜策略。

\subsection{代码实现}

\inputminted{c++}{../src/p10_not_coprime_game.cpp}


\section{总结}

本报告针对 2026 年蓝桥杯国赛 A 组的 10 道题目进行了完整的分析与求解。综观全卷，题目在算法设计层面呈现出以下几个鲜明特点。

其一，数学建模能力是解题的先决条件。题 01 将大数乘法转化为 10 的幂的线性组合从而避免高精度运算，题 02 将环形灯珠的约束转化为有向图闭游走并使用 Burnside 引理与矩阵快速幂，题 09 将平衡位置条件转化为排列不动点的组合计数——上述转换的成功与否直接决定了问题是否可解。

其二，经典算法与数据结构的灵活运用贯穿始终。题 04 通过树形 DP 在 $\mathrm{O}(n)$ 时间内解决路径计数问题，题 07 利用带权任务调度的贪心算法与权值缩放技巧实现字典序优化，题 08 在 Trie 树上运用 Grundy 数递归解决博弈问题——这些解法都是对经典思想在新的问题背景下的创造性应用。

其三，公平组合游戏的模型识别与 Grundy 数计算是两道压轴题目的核心。题 08 将字符串前缀删除游戏转化为 Green Hackenbush on a tree，题 10 则通过对因子关系的图模型分析给出了偶数和 $3$ 的奇数倍的封闭公式——体现了 Sprague--Grundy 定理在算法竞赛中的核心地位。

最后，在大规模数据约束下（$n \le 10^7$ 的排列计数、$n \le 2 \times 10^5$ 的树形 DP、$N \le 300000$ 的 Grundy 数预处理），所有题解在保证算法正确性的同时均给出了线性或近线性的实现，确保可在竞赛时限内完成运算。

\end{document}
